\documentclass[11pt,a4paper]{article}
\usepackage[T1]{fontenc}
\usepackage[utf8]{inputenc}
\usepackage{lmodern,microtype}
\usepackage[margin=25mm]{geometry}
\usepackage{amsmath,amssymb,amsthm,mathtools,mathrsfs}
\usepackage{enumitem,booktabs,array,xurl,needspace}
\usepackage{xcolor}
\usepackage[colorlinks=true,linkcolor=blue!50!black,citecolor=blue!50!black,urlcolor=blue!50!black]{hyperref}
\hypersetup{pdftitle={Weakly integral braided fusion categories below dimension 5400 squared},pdfauthor={Sebastian Burciu, Jingcheng Dong, Sebastien Palcoux}}
\allowdisplaybreaks[2]
\setlength{\emergencystretch}{2em}
\numberwithin{equation}{section}
\newtheorem{theorem}{Theorem}[section]
\newtheorem{question}[theorem]{Question}
\newtheorem{proposition}[theorem]{Proposition}
\newtheorem{lemma}[theorem]{Lemma}
\newtheorem{corollary}[theorem]{Corollary}
\newtheorem{conjecture}[theorem]{Conjecture}
\newtheorem{example}[theorem]{Example}
\theoremstyle{remark}
\newtheorem{remark}[theorem]{Remark}
\newcommand{\C}{\mathcal C}
\newcommand{\D}{\mathcal D}
\newcommand{\E}{\mathcal E}
\newcommand{\M}{\mathcal M}
\newcommand{\Pcat}{\mathcal P}
\newcommand{\Qcat}{\mathcal Q}
\newcommand{\Z}{\mathcal Z}
\newcommand{\one}{\mathbf1}
\renewcommand{\Vec}{\mathrm{Vec}}
\newcommand{\Rep}{\operatorname{Rep}}
\newcommand{\FPdim}{\operatorname{FPdim}}
\newcommand{\Irr}{\operatorname{Irr}}
\newcommand{\Hom}{\operatorname{Hom}}
\newcommand{\End}{\operatorname{End}}
\newcommand{\Inv}{\operatorname{Inv}}
\newcommand{\rank}{\operatorname{rank}}
\newcommand{\ad}{\mathrm{ad}}
\newcommand{\pt}{\mathrm{pt}}
\newcommand{\rev}{\mathrm{rev}}
\newcommand{\Gal}{\operatorname{Gal}}
\newcommand{\hs}{\widehat s}
\newcommand{\K}{K_0}
\newcommand{\kk}{\mathbb C}
\newcommand{\R}{\mathbb R}
\newcommand{\ZZ}{\mathbb Z}
\newcommand{\QQ}{\mathbb Q}
\newcommand{\CC}{\mathbb C}
\newcommand{\Tr}{\operatorname{Tr}}
\newcommand{\lcm}{\operatorname{lcm}}
\newcommand{\sVec}{\mathrm{sVec}}
\DeclareMathOperator{\ord}{ord}
\newcommand{\OO}{\mathcal O}
\title{Weakly integral braided fusion categories\\ below dimension \texorpdfstring{$5400^2$}{5400 squared}}
\author{Sebastian Burciu \and Jingcheng Dong \and S\'ebastien Palcoux}
\date{}
\begin{document}
\maketitle
\begin{abstract}
Every weakly integral braided fusion category of Frobenius--Perron dimension
strictly less than $5400^2$ is weakly group-theoretical. We reduce to two
integral character tables and use their principal blocks to restrict the prime
divisors of the dimension. For the remaining dimensions the proof uses local
character norms and short tables of residues, with all arithmetic given in the text.
\end{abstract}

\section{The bound and the reduction}
All fusion categories are over $\CC$. A category is \emph{weakly integral}
if its Frobenius--Perron dimension is an integer, and \emph{weakly
group-theoretical} (WGT) if it is Morita equivalent to a nilpotent fusion
category. The positive spherical structure identifies categorical and
Frobenius--Perron dimensions in the weakly integral case
\cite[Propositions~8.23 and~8.24]{ENOannals}.
Natale proved weak group-theoreticity below dimension $1800$ for nondegenerate
braided categories \cite[Theorem~8.1]{NataleWGT}. Dong obtained the same bound for arbitrary braided categories \cite[Corollary~3.2]{DongWeakIntegral}.

\begin{theorem}\label{short:main}
A weakly integral braided fusion category of dimension less than $5400^2=29160000$ is WGT.
\end{theorem}

The reduction is by divisibility. Given a counterexample of dimension $D_0$,
choose a braided counterexample $\C$ of least dimension among the divisors
of $D_0$, and write $M=\FPdim(\C)$. We call $M$ a \emph{minimal
obstruction}. A category or based ring is \emph{perfect} if the unit is
its only invertible simple or basis element, respectively. It is
\emph{simple} if it has no proper nontrivial fusion subcategory or based subring.

Write $\C_{\ad}$ for the subcategory generated by the simple constituents
of $X\otimes X^*$, $\C_{\pt}$ for that generated by the invertible simples,
and $\sVec$ for super vector spaces with their symmetric braiding.

\begin{lemma}\label{lem:minimal}
The category $\C$ is integral and equals its adjoint subcategory. Either it
is modular, perfect and simple, or its M\"uger center and its pointed
subcategory both equal $\sVec$. In the latter case the ring obtained by
identifying the transparent fermion with the unit is perfect and simple.
\end{lemma}
\begin{proof}
A nontrivial universal grading reduces to its neutral component, and a proper
nontrivial nondegenerate subcategory splits off with its centralizer. A
nontrivial Tannakian $\Rep(G)$ reduces to the neutral component of the
de-equivariantization, of dimension $M/(|G||H|)$ for the grading support
$H\le G$. In each reduction the smaller dimensions divide $M$, and WGT
passes back by extension, equivariantization or tensor product
\cite[Propositions~8.15 and~8.20]{ENOannals},
\cite[Proposition~4.1]{ENOsolvable}, \cite[Theorems~3.13(i) and~4.44, Proposition~4.26]{DGNO},
\cite[Proposition~4.1]{NataleWGT}. Thus $\C=\C_{\ad}$, and the
dimension grading gives integrality \cite[Theorem~3.10]{GNnil}.
Its symmetric center is $\Vec$ or $\sVec$. Every proper fusion
subcategory has dimension dividing $M$, hence is WGT; since it is integral
and contains no nontrivial Tannakian subcategory, it is pointed
\cite[Theorem~3.1]{NataleCore}. In the modular case the centralizer
of $\C_{\ad}$ is $\C_{\pt}=\Vec$. In the slightly degenerate case
$\C_{\pt}=\sVec$ \cite[Lemma~7.1]{NataleWGT}. Pulling back based
subrings from the fermion quotient gives the assertion.
\end{proof}

A \emph{minimal table} will mean the character table of one of these two
categorical reductions, with its induced ribbon and Galois data.

\begin{proof}[Proof of Theorem~\ref{short:main}]
Apply the lemma and use the ordinary table of dimension $\Delta=M$ or the
positive reduced table of dimension $\Delta=M/2$. In both cases $\Delta$
is square-full, meaning that each prime divisor has exponent at least two
(Lemma~\ref{lem:tables}). In the second case $\Delta$ is even
(Proposition~\ref{lem:principal-fermion-two}), so $M$ is square-full as well.
Dong's two-prime theorem excludes at most two prime divisors
\cite[Theorem~3.1]{DongWeakIntegral}. Six distinct primes give
$M\ge(2\cdot3\cdot5\cdot7\cdot11\cdot13)^2>5400^2$.
It remains to exclude three, four and five primes. Section~\ref{short:three-prime}
uses the quartic and quintic valuations to do so at three primes.
Sections~\ref{short:four} and~\ref{sec:short-five} use norms of local
columns to dispose of the four- and five-prime possibilities.
\end{proof}

In fact the proof only requires the square-full part
$\prod_{v_p(D_0)\ge2}p^{v_p(D_0)}$ to be below $5400^2$, since it
is divisible by $M$. For a weakly integral fusion category $\mathcal A$ the center
has dimension $\FPdim(\mathcal A)^2$ and the forgetful functor
$\Z(\mathcal A)\to\mathcal A$ is dominant
\cite[Corollary~7.13.11 and Theorem~7.16.6]{EGNO}. Thus the bound without a
braiding is $\FPdim(\mathcal A)<5400$, by descent of WGT under a quotient
\cite[Proposition~4.1]{ENOsolvable}.

\section{The character table and its residue blocks}
In the two cases of Lemma~\ref{lem:minimal} write $D=M$ for the category dimension and put respectively $\Delta=D$ and $\Delta=D/2$.
In the second case use $R=K_0(\C)/(g-1)$ for the transparent fermion $g$.
The twists give $gX\not\simeq X,X^*$, so we may take duality-compatible representatives
of the pairs. The new structure constants are $N_{XY}^Z+N_{XY}^{gZ}$.
In either case the unit coefficient $\tau$ is a perfect integral trace on $R$,
and its characters form the symmetric table $S$.
Write $d_X$ for the degree of a simple label $X$, and $f_X$ for its formal codegree. The \emph{vacuum} is the unit label $0=\one$; a row \emph{hits} a column when the corresponding entry is nonzero. The $p$-height of $X$ is $v_p(d_X)$. In both tables

\begin{equation}
 \begin{gathered}S_{\one X}=d_X,\quad S^\dagger S=\Delta I,\quad S_{X,Y^*}=\overline{S_{XY}},\\
 \chi_Y(X)=S_{XY}/d_Y,\quad f_Y=\|\chi_Y\|^2=\Delta/d_Y^2\in\ZZ.\end{gathered}
\end{equation}
For the modular table these identities are standard \cite[Sections~8.13--8.14]{EGNO}. For the reduced table, the normalized columns give all characters by \cite[Proposition~3.3]{YuSlightlyDegenerate} (see \cite[Proposition~2.7]{BGNPWSuper} for the unitary case). The ratio $\Delta/d_Y^2$ is an algebraic integer by \cite[Corollary~3.4]{YuSlightlyDegenerate}, hence an integer here. Write $I_s=\{X:d_X=s\}$, $n_s=|I_s|$ and $\sigma\chi_X=\chi_{\pi_\sigma X}$. Complex conjugation acts by duality and commutes with cyclotomic Galois action. Thus $f_{\pi_\sigma X}=\sigma f_X=f_X$, the permutation preserves degrees, and symmetry gives
\begin{equation}\label{eq:galois-permutations}
 \sigma S_{XW}=S_{\pi_\sigma X,W}=S_{X,\pi_\sigma W}.
\end{equation}
The characters take algebraic-integer values. The unit-coefficient trace is
\begin{equation}\label{eq:fusion-coefficient-trace}
 \tau(B)=\sum_X\chi_X(B)/f_X .
\end{equation}

The entries are cyclotomic \cite[Theorem~8.51]{ENOannals}. For a cyclotomic
number $z$, its \emph{house} is the maximum of the moduli of its conjugates, and $\mathsf M(z)$ is the mean of their squares.
For a nonzero algebraic integer this is at least one, by the norm.
Symmetry makes $S_{XY}$ divisible, among algebraic integers, by both $d_X$
and $d_Y$, hence by $\lcm(d_X,d_Y)$.
Balancing expresses $S_{XY}/\lcm(d_X,d_Y)$ as a nonnegative rational
combination of roots of unity of total weight $\gcd(d_X,d_Y)$:
the coefficients before division are $N_{XY}^Z d_Z$, and their sum is $d_Xd_Y$.
This uses actual representatives in the reduced case.

\begin{lemma}[Averaging over a character orbit]\label{deep:field-energy}
In ordinary or positive reduced symmetric integral data, let $H$ be a subgroup of the cyclotomic Galois group, $e=|H\cdot Z|$ and $A_H(x)=|H|^{-1}\sum_{\sigma\in H}\sigma(x)$. Then
\[
  f_Z/e=\sum_X |A_H(\chi_Z(X))|^2,\qquad
  0=\sum_X d_X A_H(\chi_Z(X))\quad (Z\ne\one).
\]
For the full Galois group write $K_X=\QQ(\chi_X)$, $F_X=\QQ(\chi_Z(X))$ and $T_X=\Tr_{F_X/\QQ}(\chi_Z(X))$. Grouping columns into orbits gives
\begin{align}
 f_Z&=\sum_{[X]}[K_Z:F_X][K_X:F_X]T_X^2,\label{eq:r8-tracenorm}\\
 0&=\sum_{[X]}d_X[K_X:F_X]T_X\qquad(Z\ne\one).\label{eq:r8-tracedim}
\end{align}
For any $H$, let $A$ be a set of nonunits whose entries $z_X=S_{XZ}/\lcm(d_X,d_Z)$ are fixed by $H$. With $\mathsf M(z)$ the mean of $|\sigma z|^2$ over all conjugates, we have
\[
 e\left(d_Z^2+\sum_{X\in A}\lcm(d_X,d_Z)^2\mathsf M(z_X)\right)\le\Delta .
\]
\end{lemma}
\begin{proof}
The $e$ conjugates of the row $\chi_Z$ are orthogonal with squared norm $f_Z$, and their average at $X$ is $A_H(\chi_Z(X))$ by \eqref{eq:galois-permutations}. Averaging preserves orthogonality to the dimension character. For the full group, the value at each of the $[K_X:\QQ]$ conjugates of $X$ is $T_X/[F_X:\QQ]$. For the inequality multiply the norm identity by $d_Z^2$, retain the vacuum and $A$, then take the full Galois mean. The Galois group is abelian, so the conjugates of these entries are also fixed by $H$.
\end{proof}

\begin{lemma}[Zeros and square-fullness]\label{lem:tables}
In a perfect table $\Delta$ is square-full. In a simple perfect table every nonunit character satisfies $|\chi_X(Y)|<d_Y$ for $Y\ne\one$.
Nonunits of relatively prime degrees have zero entry, and no nonunit has prime-power degree.
In odd dimension only the unit is self-dual.
\end{lemma}\begin{proof}
For ordinary modular categories the row-vanishing theorem is due to
Gelaki, Naidu and Nikshych \cite[Theorem~6.1]{GNNcenters}.
Square-fullness in the perfect integral modular case is
\cite[Corollary~10.10]{BP}, and the coprime-degree and prime-power
obstructions in the integral modular case follow from
Etingof, Nikshych and Ostrik \cite[Lemma~7.1 and Corollary~7.2]{ENOsolvable}.
A nowhere-zero nonunit row of degree $s$ would give
$\Delta/s^2\ge\sum_Y(d_Y/\gcd(s,d_Y))^2>\sum_Yd_Y^2/s^2$, by averaging its norm; the unit gives strictness.
Thus $\prod_X XX^*$ vanishes on all non-Perron characters. Its unit coefficient is $\prod_X d_X^2/\Delta\in\ZZ$ by \eqref{eq:fusion-coefficient-trace}. Since $d_X^2\mid\Delta$, this proves square-fullness.

The labels where $|\chi_X(Y)|=d_Y$ form a based subring by the triangle inequality on products and closure under duality. When $X\ne\one$ this subring is proper by the codegree identity. Simplicity gives strictness, so for coprime nonunit degrees every conjugate of the algebraic integer $S_{XY}/(d_Xd_Y)$ has modulus below one. Yet
\begin{equation}
 1+\sum_{Y\ne\one}d_Y\chi_X(Y)=0
\end{equation}
requires a nonunit hit of degree prime to any prescribed prime. Thus a nonunit cannot have prime-power degree.
For odd $\Delta$, all $f_X$ are odd, so $\tau(B)\equiv\Tr(N_B)\pmod2$ by \eqref{eq:fusion-coefficient-trace}. Apply this to $X^2$ and $X$, using $\Tr(N_X^2)\equiv\Tr(N_X)\pmod2$.
\end{proof}

\subsection{Character fields and local orders}
Let $K_X=\QQ(\chi_X)$ and let $e_X$ be its primitive idempotent.
Character orthogonality gives
\begin{equation}\label{eq:character-formal-codegree-general}
 f_Xe_X=\sum_Y\chi_X(Y^*)[Y].
\end{equation}
\begin{lemma}[The different and conductor]\label{cor:field-conductor}
The different of $K_X$ divides $f_X\mathcal O_{K_X}$ and its conductor divides $2f_X$.
In particular an entry common to two character fields has conductor dividing twice the gcd of their codegrees.
\end{lemma}
\begin{proof}
If $a$ belongs to the inverse different, the sum of the conjugates of $af_Xe_X$ belongs to $R$ by \eqref{eq:character-formal-codegree-general}. Its eigenvalue $af_X$ is integral.

Work after a finite unramified extension of $\QQ_p$, so the ramified
part $E$ of $K=K_X$ lies in the totally ramified cyclotomic extension
$L$ of minimal conductor $p^m$. Let $e=[E:F]$, where $F$ is the
unramified base, and let $d_E$ be the different exponent at $E$.
The bound just proved gives $d_E/e\le v_p(f_X)$.
The derivative of $\Phi_{p^m}$ gives
$d_L/[L:F]=m-1/(p-1)$.
For odd $p$ and $e=e_0p^s>1$, $e_0\mid p-1$, cyclic wild inertia
gives $m=s+1$ and $L/E$ is tame. Thus
$(d_E+1)/e=(d_L+1)/[L:F]$, or
\begin{equation}\label{eq:odd-normalized-different}
 \delta_p(e):=d_E/e=s+1-\frac{1-p^{-s}}{p-1}-\frac1e .
\end{equation}
This also gives zero at $e=1$. For $s=0$, $e>1$ it is positive,
and for $s\ge1$ it lies strictly between $s$ and $s+1$.
Hence $p^m\mid f_X$.
At two, conductor four gives normalized different one.
For $m\ge3$, minimality excludes $1+2^{m-1}$ from
$\Gal(L/E)$. Every nontrivial subgroup of $\langle5\rangle$
contains that element, so $[L:E]\le2$. The possible quadratic
generators are $-1$ and $2^{m-1}-1$. For either generator $\sigma$,
a primitive $2^m$th root $\zeta$ satisfies
$v_L(\zeta-\sigma\zeta)=2$; this is the relative different exponent.
Transitivity therefore gives
$d_E/e\ge m-1-2/2^{m-1}>m-2$, whence $2^{m-1}\mid f_X$.

The same different bound passes to subfields by the tower formula, so the conductor argument applies to common entries.
\end{proof}
The same proof applies whenever the given ring is a symmetric integral order with cyclotomic characters and this trace formula; in particular it applies to the signed order below. A character of $p$-power codegree thus lies in a $p$-power cyclotomic field
(the factor two is redundant at odd conductor) and has a unique prime above $p$.

\paragraph{Local projectors.}
Fix a prime $\mathfrak p$ of a splitting field over $p$. Write $P_p$ for the block of the dimension character. The full residue classes (or $p$-blocks) group the characters of the commutative order $R$ with equal reductions on every basis element. Over a complete unramified coefficient ring $\mathcal O$ containing the residues in its residue field, the special-fiber idempotents lift. Hence for a union $L$ of full blocks, $E_L=\sum_{X\in L}e_X$ lies in $R\otimes\mathcal O$, and for locally integral $B$,
\begin{equation}\label{eq:local}
 \tau(E_LB)=\sum_{X\in L}\chi_X(B)/f_X\in\mathcal O .
\end{equation}

\begin{lemma}[Support of a prime-power row]\label{thm:principal-support}
A row $Z$ with $f_Z=p^a$ is supported on the principal block $P_p$.
More generally, for a Galois-stable set of tests $\mathcal A$, set
$L=\{X:\chi_X(W)\equiv d_W\bmod\mathfrak p\ (W\in\mathcal A)\}$.
Then $\sum_{X\in L}|\chi_Z(X)|^2$ is an integer divisible by $p^{v_p(f_Z)}$.
\end{lemma}\begin{proof}
By \eqref{eq:galois-permutations}, permuting $X$ permutes the tests without changing their degrees, so $L$ is Galois-stable. Symmetry and \eqref{eq:local} identify the sum divided by $f_Z$ with $\tau(E_LZZ^*)$, integral at $p$. The sum is a rational algebraic integer. For all labels as tests, $L=P_p$ contains the unit, so a row of norm $p^a$ spends its entire norm there.
\end{proof}

\begin{lemma}[A local order with conductor of length one]
After an unramified base extension, write the normalization of a local symmetric block order $A$ as $M=\prod_i\mathcal O_i$ with common residue field $k$. On factor $i$ the trace is $\Tr_i/f_i$. If $e_i$ and $d_i$ denote its ramification index and different exponent, put $c_i=e_iv_p(f_i)-d_i$. A factor with $c_i=0$ is the whole block. If some $c_i=1$, there are exactly two factors, both with $c_i=1$, and $A$ is their fiber product over $k$.
\end{lemma}\begin{proof}
Self-duality gives conductor $(A:M)=M^\#=\prod_i f_i\mathfrak D_i^{-1}$; the local orders here are full lattices and their generic fibers are products of the character fields. If $c_i=0$, the factor idempotent lies in $A$. If $c_i=1$, the entire maximal ideal on this factor lies in $A$. Thus $A=\mathcal O_i\times_k B$ for the projection $B$ to the other factors, since residues must agree and all compatible pairs lift. There must be another factor: otherwise $A=\mathcal O_i$ has conductor exponent zero. Self-duality makes $A$ Gorenstein, so the fiber-product criterion makes $B$ a discrete valuation ring \cite[Fact~2.6 and Corollary~2.7]{NSTV}. The conductor is the product of the two maximal ideals.
\end{proof}

\begin{lemma}[Blocks of defect one]\label{lem:one-defect-block}
Call $v_p(f_X)$ the $p$-defect. At odd $p$ the block of a defect-one row has exactly two tame inertia orbits, both of defect one. At $p=3$ their lengths are one and two.
At two, the block of a defect-one row is either one ramified quadratic branch or a pair of branches of degrees $1,e$ and defects $1,1+v_2(e)$. Hence a block of even rank greater than two contains neither defect zero nor defect one. No unramified higher-defect row lies in a defect-one block.
\end{lemma}\begin{proof}
At odd $p$ the conductor bound makes a defect-one branch tame and $e_i-d_i=1$. An unramified defect-one branch at two has the same conductor length. The preceding lemma identifies the other branch, on which $v_p(f)=(d_E+1)/e$. At odd $p$ this quantity is one for a tame branch and strictly between $s$ and $s+1$ for wild degree $p^s>1$, by \eqref{eq:odd-normalized-different}. At three the possible lengths are $1,2$ and their sum vanishes modulo three: the block trace of the identity is $\sum e_i d_{X_i}^2/\Delta$, and after removing their common three-part the two unit squares are one modulo three.

At two a ramified defect-one branch has $2(e_i-1)\le d_i\le e_i$, by the first two ramification groups and the different bound. Thus $e_i=d_i=2$ and it is the whole block. In the two-branch case integrality of $1/f_1+e/f_2$ gives $v_2(f_2)=1+v_2(e)$. Since $e$ is a power of two, even rank $1+e$ forces $e=1$. A defect-zero branch has $e=1,c=0$ and rank one by itself.
\end{proof}

\begin{corollary}[The first principal defect]\label{cor:principal-valuation-bound}
For $v_p(\Delta)\ge2$ every principal label satisfies
$v_p(f_X)\ge2$, or equivalently $2v_p(d_X)\le v_p(\Delta)-2$.
\end{corollary}\begin{proof}
The dimension character is unramified of defect $v_p(\Delta)$. The preceding lemma excludes defect one in its block, and a defect-zero idempotent separates its own character.
\end{proof}

\begin{lemma}[Degrees detected by absolute traces]\label{prop:quartic-prime-top-codegree}
If $f_Z=p^2$, the nonunit degrees at which $\chi_Z$ has nonzero absolute trace divide $M=d_Z/\gcd(p,d_Z)$ and have gcd one. In a simple perfect table $M$ therefore involves at least three primes. In particular there is no row of codegree $p^2$ at a quartic prime of a three-prime minimal table.
\end{lemma}\begin{proof}
In \eqref{eq:r8-tracenorm} a nonzero $T_X$ at $X\ne\one$ is prime to $p$, since the vacuum already contributes. As $d_X^2\mid p^2d_Z^2$, either $d_X\mid d_Z$ or $d_X/\gcd(d_X,d_Z)=p$. The latter would make $p\mid T_X$ by lcm divisibility. If $v_p(d_X)=v_p(d_Z)>0$, then $X\in P_p$ by Lemma~\ref{thm:principal-support} and $\chi_Z(X)=(d_X/d_Z)\chi_X(Z)$ vanishes modulo every prime above $p$, again giving $p\mid T_X$.
Now \eqref{eq:r8-tracedim} expresses $-1$ as an integral combination of the indicated degrees. None is a prime power, so their gcd cannot be one on fewer than three primes. When $v_p(\Delta)=4$, $v_p(d_Z)=1$ and $M$ omits $p$.
\end{proof}

\section{Residue blocks and twist phases}
Write $s=S/\sqrt D$ and $T=\operatorname{diag}(\theta_i)$ for the ordinary modular matrices, $N=\ord(T)$, and $v_i$ (also $e_i$ in matrix calculations) for coordinate vectors. An object is \emph{$p$-active} if the $p$-primary part $\theta_{X,p}$ of its twist is nontrivial. Two primes \emph{mix} if a label is active at both; otherwise they are \emph{separated}. A prime separated from every other prime is \emph{isolated}.

For a weakly integral modular category use the lift
$W=\left(\begin{smallmatrix}0&-1\\1&0\end{smallmatrix}\right)\mapsto s$,
$U=\left(\begin{smallmatrix}1&1\\0&1\end{smallmatrix}\right)\mapsto t=T/\gamma$,
where $p^+=\sum_i d_i^2\theta_i$ and $\gamma^3=\xi=p^+/\sqrt D$ is the multiplicative central charge. We may take $\gamma$ of two-power order, equal to one for a center. The representation factors through $\mathrm{SL}_2(\mathbb Z/\ell\mathbb Z)$, $\ell=\ord(t)$ \cite[(1-10), Theorem~II(i), Proposition~6.7(ii) and Theorem~6.10]{DLNCongruence}. A prime subscript selects the local factor by the Chinese remainder isomorphism. Write $U(x)=\left(\begin{smallmatrix}1&x\\0&1\end{smallmatrix}\right)$, $L(x)=\left(\begin{smallmatrix}1&0\\x&1\end{smallmatrix}\right)$ and $H_a=\operatorname{diag}(a,a^{-1})$ for the elementary unipotents and diagonal torus.

For integral slightly degenerate $\mathcal B$ put $\Delta=\FPdim(\mathcal B)/2$ and let $g$ be its transparent fermion. With the positive spherical structure, $\mathcal B$ is pseudo-unitary, so the extension theorem of Johnson-Freyd and Reutter \cite[Corollary~1.3]{JFR} gives a minimal modular extension of dimension $4\Delta$, with matrices $s,T$. Let $Fv_i=v_{gi}$ and $Gv_i=\epsilon_i v_i$ for the fermion-monodromy sign $\epsilon_i$ (even labels have sign $+1$). The identity $S_{gi,j}=\epsilon_jS_{ij}$ shows that $s$ interchanges $F,G$ by conjugation. Since $\theta_{gi}=-\epsilon_i\theta_i$, each odd local upper unipotent commutes with both, hence so does its Weyl conjugate. These generate the odd local group. It follows that $W_2$ interchanges $F,G$ as well. On their common $+1$ space the symmetrized even vectors $(v_i+v_{gi})/\sqrt2$ form the basis of the normalized reduced matrix $\widehat s$. Even twist powers and odd local factors preserve this space. Their Bruhat and Galois identities restrict, with the same Galois sign since the two dimension normalizations differ by a rational factor.

\begin{theorem}[Block orthogonality]\label{thm:block}
For a full $p$-block $B$ of ordinary modular data,
\[
 \sum_{z\in B}\overline{s_{zi}}s_{zj}=0
       \qquad(\theta_{i,p}\ne\theta_{j,p}).
\]
In reduced integral data the same holds for $\widehat s$ at odd $p$,
and at two whenever $\theta_{i,2}^2\ne\theta_{j,2}^2$.
In particular, in the ordinary table
\begin{equation}\label{eq:zero}
 p\nmid f_z\in\ZZ,\quad p\mid\ord(\theta_j)
       \quad\Longrightarrow\quad S_{zj}=0 .
\end{equation}
\end{theorem}
\begin{proof}
First consider ordinary data. Write $L_k^i=\nu_{k,1}^{(i,i^*)}$
for the indicator with insertion in
$\Z(\C)\simeq\C\boxtimes\C^{\rev}$, extended linearly by
\cite[Corollary~5.5]{NSind}. Its modular formula on a simple $c$ is
$\sum_{a,b}s_{ia}s_{ib}(\theta_a/\theta_b)^kN_{ab}^c$
\cite[Proposition~6.1]{NSind}. Pairing with $f_ze_z$ replaces the
fusion coefficients by $\overline{\chi_z(a)\chi_z(b)}$;
duality in the second factor gives
\begin{equation}\label{eq:indicator-entry}
 L_k^i(e_z)=|(sT^ks^{-1})_{iz}|^2 .
\end{equation}

\emph{A rotation congruence.}
Put $m=N_{p'}$. Given a genuine object $A$, choose
$b\ge\max(1,v_p(N))$ with $p^b\equiv1\bmod m$, and $n\ge2b$.
Then
\begin{equation}\label{eq:large-power}
 L_m^i(A^{p^n})\equiv L_{mp^b}^i(A^{p^{n-b}})
              \pmod {p^{n-2b+1}\overline{\ZZ}_p}.
\end{equation}
Strand grouping identifies the two sides with traces of powers
$p^n,p^{n-b}$ of the same rotation on $mp^n$ strands
\cite[Lemmas~2.5, 2.7 and Proposition~2.8(ii)]{NSind}. Its
eigenvalues are $mp^n$th roots, since the insertion has twist one,
and all power traces lie in $\QQ(\zeta_N)$
by \cite[Proposition~5.7]{NSind}, since $N=\ord(T)=\operatorname{FSexp}(\C)$ by \cite[Theorem~7.7]{NSfs}. Thus Galois orbits over that field
have constant multiplicity. The two powers agree on the
prime-to-$p$ part and on roots of $p$-power order at most $p^{n-b}$.
For order $p^t>p^{n-b}$, either power is constant on the subgroup
orbit given by units $1\bmod p^b$ on the $p$-factor. Its size
$p^{t-b}$ is divisible by $p^{n-2b+1}$, proving the congruence.

\emph{Separating blocks.}
In the mod-$p$ fusion ring, residues with a given kernel and image
$\mathbb F_{p^h}$ are Frobenius conjugate. Denote their blocks
in Frobenius order by $B_0,\ldots,B_{h-1}$ (all conjugates occur
by local Galois symmetry). For $a\in\mathbb F_{p^h}^{\times}$,
lift the Chinese-remainder element with residue $a$ at this kernel
and zero at all others to a genuine object, taking nonnegative
basis coefficients. Choose $h\mid b,n$. As $n$ tends to infinity,
both powers of the object in \eqref{eq:large-power} converge in
character coordinates to $\sum_j[a]^{p^j}E_{B_j}$, where brackets
denote Teichm\"uller lifts. Hence $L_m^i$ and $L_{mp^b}^i$ agree
on that sum. The characters $a\mapsto[a]^{p^j}$ of
$\mathbb F_{p^h}^\times$ are distinct, so the functionals agree on
each $E_{B_j}$. As $N\mid mp^b$, \eqref{eq:indicator-entry} gives
\begin{equation}
 \sum_{z\in B}|(sT^ms^{-1})_{iz}|^2=\mathbf1_{i\in B}.
\end{equation}
Thus $sT^ms^{-1}$ commutes with the coordinate projector of $B$.
Conjugating back gives the assertion because $T^m$ distinguishes
the $p$-primary twists. If $p\nmid f_z$, the idempotent $e_z$
is locally integral, so $\{z\}$ is a block. Taking $i=0$ gives
\eqref{eq:zero}.

\emph{Passage to a fermion quotient.}
In its minimal modular extension select all labels whose characters
on the even subring have the prescribed residue of $B$. Their
projector $P$ is a sum of ordinary full-block projectors: equality of
residues on the full ring implies equality on the even part. Since
$z$ and $gz$ agree on the even subring, $P$ preserves the positive
fermion subspace, on which it restricts to the reduced block
projector. Put $T_p=\operatorname{diag}(\theta_{i,p})$. The ordinary
result gives $[P,sT_ps^{-1}]=0$. Restrict this identity to the
positive subspace at odd $p$, or its squared-operator version at
two: $s$ preserves that subspace and the members of a fermion pair
have equal odd primary twists and equal squared dyadic twists.
This gives precisely the two reduced assertions.
\end{proof}

\begin{lemma}[Local modular action on residue blocks]\label{lem:local-block-invariance}
Every full $p$-block subspace is invariant under the $p$-local
modular group, in ordinary weakly integral data and also in
positive reduced data at odd $p$. In particular $W_p$ preserves
each block.
\end{lemma}
\begin{proof}
Its coordinate projector commutes with $U_p$, which is diagonal,
and with $sU_ps^{-1}$ by Theorem~\ref{thm:block}, since $U_p$
differs from $T_p$ by a scalar. These elementary unipotents
generate the local factor.
\end{proof}

\begin{lemma}[Twist primes seen by a row]\label{lem:complement}
In ordinary weakly integral modular data, if the local Weyl element
$W_p$ preserves the line of $z$ then $p\nmid f_z$. In particular:
\begin{enumerate}
\item $p\mid f_z$ exactly when the $z$-row hits a $p$-active column;
\item if no twist contains both $p$ and $q$, every $p$-active
label has $q$-unit codegree.
\end{enumerate}
Every prime in the twist of $z$ divides $f_z$.
\end{lemma}
\begin{proof}
We may assume $p\mid D$ and hence $p$ divides the level by categorical
Cauchy. Choose $k$ divisible by its $p$-part and congruent to one on
the complementary level. The identity $U(WUW^{-1})U=W^{-1}$ on
the complementary factors identifies $sT^ks^{-1}$, up to diagonal
phases on the left and right, with their inverse Weyl. If $W_p$
preserves the line of $z$, the absolute value of the resulting
$(0,z)$ entry is $|s_{0z}|=1/\sqrt{f_z}$. Thus for $y=f_ze_z$,
\eqref{eq:indicator-entry} gives $L_k^0(y)=1$. Grouping $p$ strands
gives
\[
 L_k^0(y)^p\equiv L_{k/p}^0(y^p)
             =f_z^{p-1}L_{k/p}^0(y)\pmod{p\overline{\ZZ}}.
\]
Indeed the $p$th power of a trace agrees modulo $p$ with the trace
of the $p$th power. This applies first to each simple, then by
additivity and the Frobenius congruence in the commutative fusion
ring to the integral coefficients of $y$. Thus $p\nmid f_z$.

Now if the row misses all active columns, then
$sT_p s^{-1}v_z=v_z$. The local upper and lower unipotents both
preserve the line; they generate the local group, so $W_p$ also
preserves it. On this line the two unipotents have equal scalars
because they are conjugate. Comparing with $sT_p s^{-1}v_z=v_z$
gives $\theta_{z,p}=1$. Together with \eqref{eq:zero} these facts
give (1) and the twist assertion. For (2) the $q$-local group
preserves every nontrivial $p$-twist fiber. There its upper
unipotent and all its local conjugates act by scalars since
the $q$-twist is trivial. Thus $W_q$ preserves each coordinate line.
\end{proof}

\begin{lemma}[Odd active pairs]\label{lem:reduced-cauchy}
In a positive reduced table each odd $p\mid\Delta$ has an active pair, and every active pair vanishes on rows of $p$-unit codegree.
\end{lemma}\begin{proof}
If the even labels of a minimal extension were all inactive, $U_p$ and its Weyl conjugate would fix the positive fermion space pointwise: the anomaly is dyadic. The fixed space of the local group is globally invariant. It contains $v_0+v_g$, hence $T(v_0+v_g)=v_0-v_g$ and $sv_0$. Since $sv_0$ is nowhere zero, $U_p$ would fix every coordinate in the extension, contrary to categorical Cauchy \cite[Theorem~3.9]{BNRWFinite}. Vanishing follows from block orthogonality on the singleton block.
\end{proof}

\begin{lemma}[The odd twist order]\label{thm:twistdiv}
The odd part of $\ord(\theta_Z)$ divides $f_Z$, including for the reduced codegree of an even fermion pair.
\end{lemma}
\begin{proof}
If $p^k$ is the odd $p$-part of $\ord(\theta_Z)$, an automorphism fixing $K_Z$ fixes $Z$ and its square fixes $t_{ZZ}$ by modular Galois symmetry \cite[Theorem~II]{DLNCongruence}.
It therefore acts by $\pm1$ on $p^k$th roots, since the scalar normalizing $T$ is dyadic.
Thus $K_Z$ contains the real cyclotomic field at $p^k$, of conductor $p^k$ unless $p^k=3$.
The conductor bound and, for $3$, Lemma~\ref{lem:complement} give $p^k\mid f_Z$.
In a minimal extension the ordinary codegree is four times the positive reduced codegree.
\end{proof}

\begin{proposition}[Odd-valuation slices]\label{prop:odd-slice}
Suppose $p$ and $v_p(D)$ are odd. Let $B$ be a full $p$-block or the whole label set. Let $x,y$ have trivial $p$-twists, with $y$ fixed by $p$-inertia. For a complementary twist $\eta$ put
\[
 M_j=\sum_{\substack{i\in B,\ \theta_{i,p'}=\eta\\
                 \ord(\theta_{i,p})=p^j}}s_{xi}\overline{s_{yi}}.
\]
Then $M_1=(p-1)M_0$, also after summing over $\eta$. The same formula
holds in integral positive reduced data with $D$ replaced by $\Delta$,
provided the complementary fiber is taken up to the fermion sign. Thus on each such fiber the squared-dimension mass of order-$p$ twists is $p-1$ times the inactive mass, and the active proportion is at least $1-1/p$. For the entire label set, activity at all primes in any set $J$ of odd primes of odd valuation has proportion at least $\prod_{p\in J}(1-1/p)$, even conditional on fixed twist components outside $J$.
\end{proposition}
\begin{proof}
Let $E$ project onto $B\cap\{i:\theta_{i,p'}=\eta\}$. It commutes with the local $p$-factor by Lemma~\ref{lem:local-block-invariance} and commutation of distinct factors; so does $Q=sEs^{-1}$ since $s$ normalizes that factor. Write $u=T_p,w,h_t$ for its upper unipotent, Weyl and torus. For a unit $n$ modulo the local level $p^a$, the $n$th weighted moment of the $p$-twist on the fiber is
\[
 (Qwu^nw^{-1})_{xy}=(Qu^{-1/n}wh_{-n}u^{-1/n})_{xy}
 =\Big(\frac{-n}{p}\Big)(Qw)_{xy}.
\]
The middle identity is Bruhat. The outer unipotents act trivially at $x,y$. Since inertia fixes $y$, the torus acts there by the Galois sign of $d_y/\sqrt D$ \cite[Theorem~II and Remark~4.5]{DLNCongruence}; at odd $v_p(D)$ this is the displayed quadratic character, from $\sqrt{(-1)^{(p-1)/2}p}$. Average over $n$: the right side vanishes, while the left gives $M_0-M_1/(p-1)$ by the means of $p$-power roots. For the reduced formula, use the restricted projector of
Theorem~\ref{thm:block} and the complementary twists modulo sign.
The odd local group acts on the positive subspace, and its torus sign
is the same since the two dimension normalizations differ by two.
At $x=y=0$ the weights are $d_i^2/D$ or $d_i^2/\Delta$. The active bound holds on each complementary fiber and thus on every union of them; iteration over $J$ proves the product bound.
\end{proof}

\begin{corollary}[The two phases at defect one]\label{fam:M:3.9}
For odd $p$ the two orbits of an ordinary defect-one block have trivial and order-$p$ primary twists. If their lengths and degrees are $a,r$ and $b,s$, respectively, then $bs^2=(p-1)ar^2$.
\end{corollary}\begin{proof}
Here $v_p(D)$ is odd and the twist orders divide $p$ by Lemma~\ref{thm:twistdiv}. Order is constant on inertia orbits. The slice identity gives the ratio $p-1$, so exactly one of the two orbits contributes to each mass.
\end{proof}

\begin{lemma}[The visible class at defect one]\label{lem:quintic-common-class}
In an ordinary integral modular table, let $t,s$ be distinct odd primes
with $s\nmid t-1$. Suppose $\mathcal U$ is a nonempty set of
$s$-free labels of $t$-defect one and $\mathcal V$ a nonempty set of
$t$-defect-one labels in $P_s$. All these rows see only $1$ and a
single common square class of primitive $t$th roots as $t$-phases.
\end{lemma}
\begin{proof}
A defect-one block has an inactive and an active orbit; the latter
has one square class of primitive phases. Their lengths
$a_0,a_1\mid t-1$ and degrees $d_0,d_1$ satisfy
$a_1d_1^2=(t-1)a_0d_0^2$
(Lemma~\ref{lem:one-defect-block}, Corollary~\ref{fam:M:3.9}).
Thus the two degrees have equal $s$-valuation. In each block of a label
in $\mathcal U$, all active representatives have $s$-unit degree, so
every $V\in\mathcal V$ hits them by principality. Its row can see
only phases in its own block: the local Weyl preserves the block
and the complementary Weyl preserves $t$-phases. The active classes
of every such pair of blocks therefore agree.
\end{proof}

\begin{proposition}[Prime avoidance on a principal block]\label{prop:principal-support}
Let $P$ be a full principal $p$-block of ordinary integral modular data,
or of positive reduced data at an odd prime. Let $E=P$; in ordinary
data one may also intersect $P$ with the fiber of trivial
prime-to-$p$ twist. Assume that the unit is the only invertible in
$E$. Fix a label $X$ and a prime $s\ne p$. If $X$ is $p$-active,
or if $p$ and $v_p(\Delta)$ are odd and $s\nmid p-1$, then at each
prime over $s$ there is a nonunit $Y\in E$ for which
$d_Y\chi_X(Y)$ has nonzero residue. In particular $s\nmid d_Y$.
\end{proposition}
\begin{proof}
The block projector commutes with the local group
(Lemma~\ref{lem:local-block-invariance}); so does the indicated
complementary projector. Conjugation by the global Weyl normalizes
the local group. For an active $X$ the $X,0$ entry of the conjugate
projector is zero, whence $\sum_{Y\in E}d_Y\chi_X(Y)=0$.
For an inactive $X$ in the second case the slice identity
(Proposition~\ref{prop:odd-slice}) is
\[
 \sum_{\substack{Y\in E\\\ord(\theta_{Y,p})=p}}d_Y\chi_X(Y)
 =(p-1)\sum_{\substack{Y\in E\\\theta_{Y,p}=1}}d_Y\chi_X(Y).
\]
In either identity the vacuum contribution is a unit at $s$, so a
nonunit term must survive modulo each prime over $s$.
\end{proof}

In either setting of the proposition, at valuation two or three the
nonunit degrees in $P_p$ of a perfect table are $p$-free and have
gcd one: apply prime avoidance to an active column at every other prime.

\subsection{The fermion at two}
We use the full modular extension to recover the information that squaring the twists loses.

\begin{lemma}[An integral minimal modular extension]\label{lem:integral-minimal-extension}
Let $\mathcal B$ be integral and slightly degenerate, with
$\dim\mathcal B=D>2$, pointed subcategory $\sVec=\langle g\rangle$,
and no other proper nontrivial fusion subcategory.
Every minimal modular extension has pointed subcategory $\sVec$,
adjoint subcategory $\mathcal B$, and precisely the proper fusion
subcategories $\Vec,\sVec,\mathcal B$.
Moreover one can choose an integral minimal modular extension.
In such an extension, $|S_{XY}|<d_Xd_Y$ for noninvertibles $X,Y$; if their degrees are coprime then $S_{XY}=0$. At odd $p$, a positive $p$-principal pair has a full $p$-principal representative. When at least three primes divide $D$, no noninvertible of the extension has prime-power degree.
\end{lemma}
\begin{proof}
By the extension theorem of Johnson-Freyd and Reutter \cite[Corollary~1.3]{JFR} there is a minimal modular extension
$\M_0$, of dimension $2D$, with neutral component $\mathcal B$
for fermion monodromy. An odd invertible $h$ would give a pointed
subcategory of order four whose bicharacter has trivial radical:
$b(g,g)=1$, $b(g,h)=-1$, and in the cyclic case $b(h,h)^2=-1$.
Splitting it off gives $\mathcal B=\sVec\boxtimes\mathcal N$ with
$\dim\mathcal N=D/2>1$, impossible. Hence
$(\M_0)_{\pt}=\sVec$ and $(\M_0)_{\ad}=\mathcal B$.
A proper subcategory crossing the grading has twice the dimension
of its intersection with $\mathcal B$, hence dimension two or four.
Without an odd invertible, its dimension must be four: its even part has $1,g$ and its odd mass two is a single label of dimension $\sqrt2$. Rigidity gives the Ising fusion rules, so this braided subcategory is nondegenerate by \cite[Corollary~B.12]{DGNO}. It splits off with fermion $g$, giving the same contradiction.

Normalize $t_X=\theta_X/\zeta$ by choosing a two-power root
$\zeta$ with $\zeta^3=p^+/\sqrt{2D}$ \cite[Theorem~6.10]{DLNCongruence}.
At an odd prime any inertia element $\sigma$ fixes $t_0$.
The Galois image of the unit is invertible and its normalized twist
is $\sigma^2(t_0)=t_0$; thus it is the unit, since $\theta_g=-1$
\cite[Theorem~II]{DLNCongruence}. The vacuum-column ratios $d_X$
are consequently unramified at every odd prime. Since $d_X^2$ is an
integer, $d_X$ is an integer or $\sqrt2$ times one.

In the nonintegral case the integral subcategory is exactly
$\mathcal B$. Let $\mathcal I$ be Ising with fermion $\psi$ and
noninvertible $\sigma$. Condense the diagonal boson
$g\boxtimes\psi$ in $\M_0\boxtimes\mathcal I$: the local modules
have dimension $(2D)4/2^2=2D$ and are nondegenerate
\cite[Propositions~4.26, 4.55(i) and~4.56(i)--(ii)]{DGNO}. Free modules on $X\boxtimes\one$,
$X\in\mathcal B$, embed $\mathcal B$ braided fully faithfully since
the nontrivial algebra summand contributes no morphisms by adjunction.
These give all even pairs after identifying boson translates.
An odd pair $(Y,\sigma)$ has dimension $2n$; a free orbit gives
one simple of dimension $2n$, and a fixed orbit gives two of dimension
$n$ since a cyclic stabilizer has only one-dimensional projective
irreducibles. This extension is integral.

For noninvertible $X$ in an integral extension,
$\langle XX^*\rangle=\mathcal B$: containment in $\sVec$ gives
$d_X^2\le2$ since $1,g$ each occur at most once.
Equality $|S_{XY}|=d_Xd_Y$ implies projective centrality, hence $Y$ centralizes $XX^*$ and belongs to $\mathcal B'=\sVec$ \cite[Proposition~8.22.5 and Lemma~9.15.1]{EGNO}. Thus there is strict inequality; for coprime degrees the lcm-normalized integral entry has all conjugates in the open unit disk and vanishes.

Take a representative of a positive $p$-principal pair, $p$ odd, and write $\chi$ for its full character. Its residue agrees with dimension on even objects. For all odd $A,B$ the product $AB$ is even, so $\chi(A)\chi(B)\equiv d_A d_B$. If an odd $A$ has $p$-unit degree, its ratio is $\pm1$ by taking $B=A$ and the same sign works throughout the odd sector. If none does, squaring gives zero residues throughout it. Tensoring the representative by $g$ reverses the sign on the odd sector. Finally, given a putative noninvertible of degree $t^i$, take an odd dimension prime $p\ne t$ and a $p$-active even pair by Lemma~\ref{lem:reduced-cauchy}. Prime avoidance on the positive block $P_p$ gives a $t$-free nonunit pair there. Its full principal representative must hit the putative degree, whose residue at $p$ is nonzero, contradicting coprime vanishing.
\end{proof}

\begin{lemma}[The two fermion quotients]\label{lem:fermion-quotients}
In an integral extension $\M$ of $\mathcal B=\langle g\rangle'$, let $R_\pm=K_0(\mathcal B)/(g\mp1)$. These are symmetric orders with identical special fibers at two. The restrictions of modular characters to the even sector have exactly the fermion orbits as fibers; even labels yield the positive characters and odd labels the signed characters. Their codegrees are
\begin{equation}\label{eq:fermion-restriction-codegree}
 f_X^\pm=\begin{cases}f_X/2&gX=X,\\f_X/4&gX\ne X,\end{cases}
 \qquad f_X=\dim(\M)/d_X^2 .
\end{equation}
Dimension residue on the even sector is equivalent to dimension residue on the full extension. The full principal two-block $P$ satisfies
\begin{equation}\label{eq:principal-fermion-masses}
 \sum_{X\in P\cap\mathcal B}d_X^2=\sum_{X\in P\setminus\mathcal B}d_X^2 .
\end{equation}
If $\Delta=2^s m$ with $s=2,3$ and $m$ odd, the positive principal degrees are odd; the full principal degrees have two-valuation at most one. Both quotient principal blocks have the same rank $k$, divisible by $2^s$.
\end{lemma}\begin{proof}
In either quotient, the trace on duality-compatible pair representatives satisfies $\tau(XY^*)=\delta_{XY}$. Reduction at two identifies the orders and hence the ranks of corresponding blocks. In the full table $\chi_{gX}(Z)=\epsilon_Z\chi_X(Z)$ for the monodromy sign $\epsilon_Z$. Orthogonality with the even projector gives
\begin{equation}
 \sum_{Z\in\mathcal B}\chi_X(Z)\overline{\chi_Y(Z)}
 =\tfrac12 f_X(\delta_{XY}+\delta_{X,gY}).
\end{equation}
This determines the fibers. Every even-subalgebra character extends by semisimplicity; its value at $g$ selects the quotient. The norm over representatives is half the even sum, giving \eqref{eq:fermion-restriction-codegree}. Its rational value is integral as a Casimir eigenvalue.
For odd $Y$, $Y^2$ is even and squaring is injective in the residue field, so dimension residue extends. Since $1,g\in P$ have opposite dyadic twists, block orthogonality gives \eqref{eq:principal-fermion-masses}. Apply Corollary~\ref{cor:principal-valuation-bound} at dimensions $\Delta,4\Delta$ for the degree bounds. In the positive principal block $2^s$ divides the sum of the odd squares $d_X^2$ by its identity trace. Each square is one modulo $2^s$.
\end{proof}

\begin{lemma}[Unramified signed characters from sign-only twists]\label{lem:signed-unramified-sign-twists}
Let $\M$ be an integral modular category with fermion $g$, and let
$\mathcal B=\langle g\rangle'$. If the two-primary part of every
actual twist of $\mathcal B$ is $1$ or $-1$, then every character
field of $\K(\mathcal B)$ is unramified at two. In particular this
holds for the signed quotient $\K(\mathcal B)/(g+1)$.
\end{lemma}
\begin{proof}
Let $\sigma$ belong to inertia at two. The Galois identities
\cite[Theorem~II]{DLNCongruence} give
$\sigma\chi_Y=\chi_{\widehat\sigma(Y)}$ and
$t_{\widehat\sigma(Y)}=\sigma^2(t_Y)$, $t_Y=\theta_Y/\zeta$.
Since the dimension character is rational, $\widehat\sigma(0)=0$
and $\sigma^2(\zeta)=\zeta$. For even $Y,Z$ the balancing formula
expresses $S_{YZ}$ in integral dimensions and actual even-sector
twists, all in an odd cyclotomic field. Thus $Y$ and
$\widehat\sigma(Y)$ have the same restricted character and lie
in one fermion orbit by Lemma~\ref{lem:fermion-quotients}.
Their normalized twists agree, while $Y,gY$ have opposite twists.
Hence $\widehat\sigma(Y)=Y$. Symmetry of $S$ and rationality of
dimensions now give $\sigma(\chi_X(Y))=\chi_X(Y)$ on all even $Y$
for every $X$. Every character of the even subalgebra extends.
\end{proof}

\begin{proposition}[Odd supports in the principal two-block]\label{lem:principal-fermion-two}
Choose an integral minimal extension $\M$ of a slightly degenerate
minimal counterexample $\mathcal B$. Its positive reduced dimension is even. Two odd primes $p,q$ of
valuation two in $\dim\mathcal B$ cannot occur as a simple degree $pq$.

If all principal two-block degrees are odd, their noninvertible
supports intersect pairwise with empty total intersection. With
exactly three odd primes each of their three pairs must occur as a
support. When $v_2(\dim\mathcal B)=3$ all principal degrees are
indeed odd. Thus this two-valuation is excluded if the odd part has
at most two prime divisors, or three with two valuations equal to two.
\end{proposition}
\begin{proof}
If $\Delta=\dim\mathcal B/2$ were odd, the positive dimension
idempotent would be integral at two. Thus the full principal block
would have only $1,g$ in its even sector by Lemma~\ref{lem:fermion-quotients}.
Its odd mass would be two by \eqref{eq:principal-fermion-masses},
forcing an odd invertible, contrary to Lemma~\ref{lem:integral-minimal-extension}.

A label of degree $pq$ has odd residue, hence every noninvertible
principal two-block label meets $\{p,q\}$ in its support by coprime
vanishing. Choose even $p$- and $q$-active columns $Y,Z$ by reduced
Cauchy. A $p$-bearing label has $p$-unit codegree and vanishes at
$Y$; likewise at $q,Z$. Thus $1,g$ are their only common principal
hits. Replacing $Z$ by $gZ$ makes the dyadic phases distinct while
both contributions stay positive, contrary to Theorem~\ref{thm:block}.

If all principal degrees are odd, any two noninvertible labels
meet by principality and coprime vanishing. They have no common
prime: modulo that prime the mass identity
\eqref{eq:principal-fermion-masses} would read $2=0$, from the
even invertibles $1,g$ alone. This also excludes an empty family.
On two primes intersecting supports have a common prime; on three,
empty total intersection forces all three pairs. The squared pair
would then occur as a degree, contrary to the first assertion.

If $v_2(\dim\mathcal B)=3$, Lemma~\ref{lem:fermion-quotients}
gives signed principal rank divisible by four. An even principal label
would be odd-sector of degree twice odd and signed defect zero or one
by \eqref{eq:fermion-restriction-codegree}, contradicting
Lemma~\ref{lem:one-defect-block}.
\end{proof}

\subsection{Two principal residues}
The \emph{weight} of a cyclotomic integer is at most $g$ when it is a nonnegative rational combination of roots of unity of total coefficient at most $g$.

\begin{proposition}[Ramification detected by a residue]\label{deep:positive-inertia}
Suppose every conjugate of a cyclotomic integer $z$ has modulus below $g$ and $p\ge g$.
If $z\ne0$ has residue zero at every prime above $p$, then $p$ ramifies in $\QQ(z)$.
The same conclusion holds if $z$ has residue $g$ above $p$ and above a second prime $q$.
\end{proposition}\begin{proof}
In an unramified field the first residue would give an algebraic integer $z/p$ of house below one, hence zero. In the second case $1+(z-g)/p$ would likewise be integral of house below one, forcing $z=g-p$, contrary to the residue at $q$.
\end{proof}
\begin{theorem}[Short sums of roots of unity]\label{three:word}
Let $z$ have weight at most $g\in\ZZ_{>0}$ and residue $g$ at every prime above a rational prime $r>g$.
If $r>2g$, it is a unique unsigned sum of $g$ roots of $r$-power order (at any common level).
Suppose also that $z$ has residue $g$ above a second prime $q<r$. If $g\le q$ then $|z|=g$.
For $(g,q,r)=(4,3,5)$ and $|z|<4$, tame five-inertia has orbit size four on $z$.
\end{theorem}\begin{proof}
At level $s^a$ the relations on powers of $\zeta_{s^a}$ are constant on each coset of $\mu_s$, by $\Phi_{s^a}$. This is the prime-power case of the R\'edei--de Bruijn--Schoenberg theorem \cite[Theorem~5.1]{Malikiosis}. This remains true modulo a different prime. In particular an unsigned word of length $<s$ is unique: on each coset its minimum multiplicity is zero.

First take a signed integral word with at most $g$ terms on the indicated primes. For just $r>2g$, erasing the $r$-parts gives a signed count in $[-g,g]$ congruent to $g$ modulo $r$. All $g$ signs are positive. For two primes let $w$ be the signed sum after erasing the $r$-parts. It lies in a $q$-power field, unramified at $r$, so $(w-g)/r$ is integral. Every conjugate of $1+(w-g)/r$ lies in the unit disk and meets the boundary only at $1$. Its norm forces $w=g$ or $w=g-r$.

Instead erase the $q$-parts, obtaining signed multiplicities $c_i$ on the $r$-power roots, with position $0$ representing $1$. The coset relations say $c_i-g\delta_{i0}$ is constant modulo $q$ on each $\mu_r$-coset. If each coset has an unused position even after marking $0$, all constants vanish. The signed count is then $g$ modulo $q$ and equals the residue of $w$ above $q$. This rules out $w=g-r$. As $w=g$, all $g$ erased summands equal $1$. Thus the original word is unsigned and $r$-primary, with every nonidentity multiplicity divisible by $q$ (for $q=2$ signs do not affect parity). If $g\le q$ all its terms coincide. At $(4,3,5)$ a strict word is $1+3\xi$ with $\xi\ne1$ of five-power order, whose tame orbit has size four.

If a coset has no unused position, $r>g$ forces $g=r-1$ and one position on each nonidentity $r$th root. All signs are $-g$ modulo $q$; a positive sign would give $q\mid r$, impossible. Thus $q\mid g-1$, excluding $g\le q$. At $(4,3,5)$ the word is $-\sum_{t=1}^4\eta_t\zeta_5^t$ with $\eta_t$ of three-power order. Here $w=-1$ since the negative summands cannot sum to four, so $\sum\eta_t=1$. The four indicated fifth roots are linearly independent over the three-power field. A tame stabilizer acts freely on their indices and preserves their coefficients. Its order divides the algebraic integer $1$, hence it is trivial.

To pass from rational weights to signed integral words, split a common level as $MN$, with $M$ supported on the one or two retained primes and $(M,N)=1$. The signed $M$th roots form a unimodular configuration \cite[Theorems~3, 13]{BH}, so the minimum weight of an integral point is attained by an integral word. For any $N$th root $a\ne\pm1$, the lattice $\ZZ[\zeta_N]$ has an integral functional fixing $1$, killing $a$ and taking roots to $0,\pm1$: in each primary factor omit one root per prime-order coset for a basis, retaining $1$ and, in a factor where the component of $a$ is not a sign, that component. Take the coefficient of $1$ in the tensor basis. Contract along this functional, using $\ZZ[\zeta_{MN}]=\ZZ[\zeta_M]\otimes\ZZ[\zeta_N]$. Each retained residue is preserved since the radical at $s\mid M$ is generated by $1-\zeta_{s^{v_s(M)}}$. The contracted point has minimum weight at least $g$ by the integral calculation above. Killing a positive-weight summand would give weight strictly below $g$, so every complementary component is a sign. Unimodularity now applies to $z$ itself.
\end{proof}

\begin{lemma}[Complementary principal degrees at three primes]\label{three:complementary}
Let a perfect fusion-simple ordinary or positive reduced table have
three dimension primes $t,u,v$. In the reduced case assume $u$ odd.
The principal block $P_u$ contains $v$-free nonunits. Every such
degree is $t^iu^j$ with $i,j\ge1$ and
$2j+2\le v_u(\Delta)$. Thus at a quartic $u$ both complementary
degrees $ut^i,uv^k$ occur in $P_u$, each of defect two.

When $u<v$ are odd, the $t$-parts of all $v$-free nonunits in $P_u$
and all $u$-free nonunits in $P_v$ are at least $u$.
\end{lemma}
\begin{proof}
Prime avoidance (Proposition~\ref{prop:principal-support}) for a
$u$-active column produces a $v$-free nonunit. It cannot have
prime-power degree, and Corollary~\ref{cor:principal-valuation-bound}
bounds its $u$-height. At $v_u(\Delta)=4$ interchange $t,v$.
For $u<v$ odd pair opposite labels $X,Y$ of degrees
$t^iu^j,t^kv^l$. The lcm-normalized entry has weight
$g=t^{\min(i,k)}$ by balancing and strict house below $g$.
Principality gives residue $g$ over both $u,v$, on all conjugates
by Galois permutation of labels and symmetry. Theorem~\ref{three:word}
therefore gives $g\ge u$.
\end{proof}

\begin{corollary}[The smallest exponent at three primes]\label{prop:reduced-quartic-two-exclusion}
If a minimal obstruction has three repeated primes in its source dimension, each has exponent at least four. In a reduced table with two odd primes $q<r$ one has $2^{\lfloor v_2(\Delta)/2\rfloor}\ge q$.
\end{corollary}
\begin{proof}
Every dimension prime of the ordinary or positive reduced table is repeated by square-fullness. At least three remain by Dong's two-prime theorem. Lemma~\ref{three:complementary} gives the valuation bound at every ordinary prime and every odd reduced prime. In the reduced case Proposition~\ref{lem:principal-fermion-two} gives $2\mid\Delta$, and the two odd principal blocks have a strictly nonzero crossing with common two-part at least $q$ by Lemma~\ref{three:complementary}. Square-divisibility bounds that two-part above as asserted.
\end{proof}

\section{Local Weyl columns}
Throughout this section $D=n^2$ is the ordinary modular dimension, except in the last ramification assertion of Proposition~\ref{deep:cluster}.

\begin{lemma}[Local Weyl vacuum column]\label{lem:local-vacuum-column}
Let $\C$ be integral modular of dimension $D=n^2$ with positive
spherical structure and normalized Weyl matrix $s=S/n$. Use the
linear lift $t=T/\gamma$ of Lemma~\ref{lem:complement}, with
$\gamma^3=\xi=p^+/n$ and $\gamma$ of two-power order
($\gamma=1$ for a center). For a set $J$ of prime divisors of the
level $\ord(t)$ put $W_J=\prod_{p\in J}W_p$. Choose a positive $m$
congruent to one at the level factors in $J$ and zero at the others.
Then $w=W_Je_0=W_J^{-1}e_0$ is rational, $nw$ is integral, and
\[
 A_J:=n^{-2}\sum_i d_i^2\theta_i^m=\xi^m (W_J)_{00}.
\]
For $q\mid n$ outside $J$, $q\mid(nw_i)^2$ for every $i$.
In particular $w_i$ is $q$-integral if $v_q(n)=1$.
For $J=\{p\}$ its support lies in the principal $p$-block.
\end{lemma}
\begin{proof}
The Galois identity
$\sigma_k(\rho(g))=\rho(A_k g A_k^{-1})$,
$A_k=\operatorname{diag}(k,1)$, follows on the two generators from
\cite[Theorem~II and Remark~4.5]{DLNCongruence}. The torus permutations
have positive signs and fix the vacuum because $s_{i0}=d_i/n$ are
positive rational numbers. Conjugating $W_J$ introduces only a torus
factor on its right, proving rationality. The local torus at $-1$
fixes $e_0$ as well, giving the inverse Weyl assertion.
In the factors in $J$ apply $U(WUW^{-1})U=W^{-1}$. Its vacuum
entry gives the displayed identity since the three scalar factors
contribute $\gamma^{-3m}$. It also identifies the absolute entries
of $sT^m s^{-1}e_0$ with those of $w$. The regular object
$R=\sum_i d_i[i]=n^2 e_{\one}$ thus has indicator
$L_m^i(R)=n^2|w_i|^2$ by \eqref{eq:indicator-entry}. This is both a
rational number and an algebraic integer, so $nw_i\in\ZZ$.
For $q\mid n$ outside $J$ choose $m$ also divisible by $q$. Grouping
$q$ adjacent strands, as in Lemma~\ref{lem:complement} gives
$L_m^i(R)^q\equiv L_{m/q}^i(R^q)\equiv0$ above $q$:
the insertion $(i,i^*)$ has twist one, and
$R^q=(n^2)^{q-1}R$. Therefore $q\mid(nw_i)^2$ in $\ZZ$.
Finally, the projection onto the principal $p$-block commutes with the
$p$-local factor by Lemma~\ref{lem:local-block-invariance}, so
$W_pe_0$ belongs to that block.
\end{proof}

\begin{proposition}[Weyl columns on separated prime sets]\label{deep:cluster}
Partition the primes in the modular level of an integral modular category
of dimension $n^2$ into sets $J$ such that each simple has twist
supported in at most one set. Write $\gamma_J$ for the $J$-part
of the anomaly normalization $\gamma$, $n_J$ for the $J$-part of
$n$, and $V_J=\gamma_J^{-3}W_J$. Then
\begin{equation}\label{eq:weylsplit}
   (V_I-1)(V_J-1)=0\quad(I\ne J),\qquad
   s=\xi\Bigl(1+\sum_J(V_J-1)\Bigr),\qquad
   n_J W_J e_0\in\mathbb Z^{\Irr(\C)} .
\end{equation}
If $p\in J$ and $v_p(d_X)<v_p(n)$, then
$v_p((n_J W_Je_0)_X)=v_p(d_X)$. If every other prime of $J$
occurs to full height in $d_X$, then $X\in P_p$.

A character field ramifies at primes of at most one set $J$. More
generally, in any integral modular category a character field cannot
ramify at two distinct odd dimension primes that never occur together
in a twist.
\end{proposition}
\begin{proof}
In factor $J$ let $U_J$ denote the represented upper unipotent. Its
normalized twist $\gamma_J U_J$ has the actual $J$-phases. Let $E_J$
be the span of the translates under the $J$-factor of the nontrivial
phase spaces. On each of these spaces the complementary upper
unipotents are scalar; their generating conjugates have the same
scalars since the entire complementary group preserves the space.
The same holds on its translates, by commutation. Hence on $E_J$
each $W_I$, $I\ne J$, is the scalar $\gamma_I^3$ by
$U(WUW^{-1})U=W^{-1}$. On $E_J^\perp$ the $J$-upper unipotent is
scalar $\gamma_J^{-1}$, so the same reasoning gives $W_J=\gamma_J^3$.
This proves the zero products and the formula for
$s=\prod_J W_J$. In particular on $E_J$ one has
$s=(\xi/\gamma_J^3)W_J$.

Choose $m$ selecting $J$ as in Lemma~\ref{lem:local-vacuum-column}.
Up to diagonal phases and duality, the vacuum coordinates of $W_J$
are the entries of $sT^m s^{-1}e_0$. Subtracting the untwisted
coordinate leaves a sum of $(\theta_Y^m-1)\chi_Y(X)/f_Y$.
Only $J$-active labels occur; by Lemma~\ref{lem:complement} their
codegrees are supported in $J$. Thus $W_Je_0$ is integral away from
$J$, while $nW_Je_0$ is integral by Lemma~\ref{lem:local-vacuum-column}.
In the split vacuum column all terms except the $J$-term are
$p$-integral, proving the valuation assertion by comparison with
$d_X/n$. Full height at $q\in J\setminus\{p\}$ gives $q\nmid f_X$;
the proof of Lemma~\ref{lem:complement} shows that $W_q$ preserves
the $X$-line. The nonzero $X$-entry of $W_J e_0$ therefore gives a
nonzero entry of $W_p e_0$ there, proving principality.

For the ramification assertion take two distinct primes in different
sets, or any two separated odd primes. The same moving-space argument
gives $(P-I)(Q-I)=0$ for their inertia torus permutations: the
complementary torus acts as a scalar on an active phase space,
hence as the identity there because it permutes its coordinate vectors.
Here the upper unipotent at the odd prime acts trivially on the
complement of its moving space, and so does its torus. The torus signs
are positive in square dimension. In arbitrary
integral dimension both separated odd primes have even valuation by
Lemma~\ref{lem:complement} applied to the opposite active sectors, so
their local Galois actions also fix $\sqrt D$ and have positive sign.
If both permutations moved the label $X$, the diagonal entry
at $X$ of $(P-I)(Q-I)$ would be positive.
\end{proof}

\begin{proposition}[Simultaneous singleton identities]\label{second:two-columns}
Call $t$ \emph{isolated} if every $t$-active label has pure $t$-power
twist, and a \emph{singleton} if $v_t(n)=1$. Let $T$ be a nonempty
set of odd isolated primes, each occurring
once in $n$. Write $m_t=n/t$, $\epsilon=\xi$,
$j_t=t(1-(W_t)_{00})$ and $a_t(X)=\epsilon t(W_t)_{X0}$ for $X\ne0$.
Then $\epsilon\in\{1,-1\}$, $j_t>0$ (and $j_t<t$ if $m_t>1$),
$m_tj_t\equiv-\epsilon\pmod t$, and
$\sum_{t\in T}j_t/t\le1+1/n$.
Every $t$-active label has degree $m_t$, and their number is at least
$(t-1)j_t$. One has
\begin{equation}\label{eq:isolated-deficit}
 tj_t=\sum_{Y\ t\text{-active}}(1-\theta_{Y,t}).
\end{equation}
For odd $n$, each $j_t$ is even. Each $a_t$ is integral, supported
exactly on the nonunit $t$-free degrees, which belong to $P_t$. The
column identities are
\begin{align}
 a_t(X)&\equiv d_X/m_t\pmod t,&
 \sum_{X\ne0} d_Xa_t(X)&=(n+\epsilon)j_t,\nonumber\\
 \sum_{X\ne0}a_t(X)^2&=2tj_t-j_t^2,&
 \sum_{X\ne0}a_u(X)a_v(X)&=-j_uj_v\quad(u\ne v).\label{eq:budget}
\end{align}
In particular, after the active labels are removed the residual
squared norm is at most $B_t=j_t(t+1-j_t)$.
\end{proposition}
\begin{proof}
Use the singleton clusters $\{t\}$, $t\in T$, and their complement $H$. A $t$-active label $Y$ has
$f_Y=t^2$, by complementarity and $v_t(f_Y)=2-2v_t(d_Y)>0$.
The equality $s=\xi W_t$ on $E_t$ gives $(W_t)_{Y0}=1/(t\xi)$;
rationality gives the common sign. The vacuum Gauss formula gives
\eqref{eq:isolated-deficit}, so $j_t>0$. The mean of a nontrivial
$t$-power root is at least $-1/(t-1)$, proving the count. For odd
$n$ the active labels come in distinct equal-twist dual pairs, and
$t$ is odd.

Write $z_t=(W_t-I)e_0$. By \eqref{eq:weylsplit},
$s=W_H+\epsilon\sum_{t\in T}(W_t-I)$, with $W_H=I$ when $H$ is
empty and $\epsilon=1$ in that case. Unitarity at the vacuum gives
the deficit-sum bound. Multiply the split vacuum column by $n$ and
reduce modulo $t$ to get both congruences. In particular $j_t\ne t$,
so $j_t<t$ when $m_t>1$. The scaled column is integral and its vacuum
entry is nonzero, so every nonunit coordinate has absolute value $<t$. This proves the
support and principality assertions. Finally $z_u\perp z_v$ for
$u\ne v$, $\|z_t\|^2=2j_t/t$, and
$\langle z_t,se_0\rangle=\epsilon\langle z_t,W_te_0\rangle
=\epsilon j_t/t$. Subtract the vacuum in these three identities to
obtain \eqref{eq:budget}; each active coordinate has square one.
\end{proof}

\begin{lemma}[Energy of a prescribed residue]\label{lem:singleton-residue-energy}
If $p$ is odd and an algebraic integer $x$ in a $p$-power cyclotomic field has residue $k$ at $p$, $0\le k<p$, then
$\mathsf M(x)\ge k(p-k)/(p-1)$.
\end{lemma}
\begin{proof}
At conductor $p$ choose coefficients $c_i\in\ZZ$ in the $p$th roots with $\sum c_i=k$
by adding a common integer. Averaging their squared sum gives
$(p\sum c_i^2-k^2)/(p-1)$, and $\sum c_i^2\ge\sum c_i=k$.
At conductor $p^a$ expand in the integral basis
$1,\zeta,\ldots,\zeta^{p^{a-1}-1}$ over $\QQ(\zeta_p)$, $\zeta=\zeta_{p^a}$.
These powers are orthogonal for the relative trace inner product, so their coefficient means add.
The coefficient residues add to $k$ modulo $p$.
Writing $h(u)=u(p-u)$ for residues in $[0,p)$, one has $h(u)+h(v)\ge h(u+v\bmod p)$, proving the assertion.
\end{proof}

\begin{lemma}[Column norm on isolated row sets]\label{lem:singleton-energy}
In the setting of Proposition~\ref{second:two-columns}, let
$X$ be any nonunit label of degree $d$ and $[d]_t$ its centered
residue at $t$. For $k=|[d]_t|$ put $E_t=j_tk(t-k)/t^2$.
Then $\sum_{t\in T}E_t<1$. If $K_X$ is unramified at $t$ one may replace
$E_t$ by $(t-1)j_t[d]_t^2/t^2$. If also $t\ge5$, every $t$-active
$Z$ has $\chi_Z(X)=[d]_t$. For such an inactive $X$,
\[
   (W_t-I)e_X=[d]_t(W_t-I)e_0,\qquad a_t(X)=-\epsilon j_t[d]_t .
\]
Any set $L$ of nonunit inactive $t$-unramified labels therefore satisfies
$1+\sum_{X\in L}[d_X]_t^2\le\lfloor(t+1)/j_t\rfloor$.
In particular, for every isolated odd singleton, $j_t\mid t-1$ or
$j_t\mid t+1$.
\end{lemma}
\begin{proof}
Each active $Z$ is $t$-principal, of codegree $t^2$. Its value on
$X$ lies in a $t$-power cyclotomic field and reduces to $d$ at $t$.
Lemma~\ref{lem:singleton-residue-energy} and the active count give
at least $E_t$ in the Galois mean of the normalized $X$-column norm.
Distinct active sectors are disjoint and the vacuum also contributes.
Apply the estimate at the active degree $m_t=n/t$, with
$k=|[m_t]_t|$. Then $0<j_tk<2t^2/(t+1)<2t-1$ and
$j_tk\equiv\pm1\pmod t$, so $j_tk$ is $1$, $t-1$ or $t+1$.

If $K_X$ is unramified at $t$, symmetry gives rational integers
congruent to $d$, hence the sharper contribution. Distinct integers congruent modulo $t$ have squares summing
to at least $(t^2+1)/2$. Each occurs on an active orbit of at least
$(t-1)/2$ labels by twist covariance, exceeding the norm budget for
$t\ge5$. The one value $c$ obeys $(t-1)j_tc^2<t^2$ and hence is
$[d]_t$. If $X$ is inactive, $e_X-ce_0$ and its inverse Weyl image
both vanish on the active sector, so the upper and lower unipotents
fix it. The vacuum entry gives the formula for $a_t$ since
$W_t^{-1}e_0=W_te_0$. On the active rows, columns $0,L$ now form
a rank-one matrix with constant row $t^{-1}(1,([d_X]_t)_{X\in L})$.
Its operator norm is at most one. The active count and integrality
of $1+\sum[d_X]_t^2$ give the stated bound.
\end{proof}

\begin{lemma}[Prime supports at an isolated prime]\label{lem:isolated-dyadic-square}
Let $t$ be isolated in an ordinary minimal table of square dimension.
Suppose a nonunit of codegree $t^2$ has $t$-valuation at most one in
its degree. The proper nonunit $t$-free degrees intersect pairwise in
prime support, have gcd one, and include at least three distinct degrees.
Here proper means smaller than the $t$-free part of $n$.

In particular, every isolated singleton $t$ satisfies $t\ge5$.
Writing $m=n/t$, its nonunit $t$-free labels satisfy
\[
 \sum_X\left(a_t(X)^2-\frac{d_X}{m}a_t(X)\right)
       =j_t(t-j_t)-\frac{\epsilon j_t}{m}.
\]
\end{lemma}
\begin{proof}
Lemma~\ref{prop:quartic-prime-top-codegree} gives $t$-free
nonunit degrees of gcd one. They divide the $t$-free part of $n$,
so deleting that full degree preserves the gcd.
All $t$-free labels are in $P_t$ by Proposition~\ref{deep:cluster}.
Two such nonunits hit one another by principality, so share a prime.
Pairwise intersection with empty total intersection requires at least
three distinct proper degrees and at least three distinct primes in them.

At a singleton a $t$-active row exists and has degree $m=n/t$ and
codegree $t^2$ by complementarity.
The active row hits the vacuum, the full degree and all three
proper degrees by principality. Each costs at least one in its
Galois-mean squared norm, excluding $t=2$.
For odd $t$ subtract the dimension trace divided by $m$ from the
squared norm in \eqref{eq:budget}. Each displayed summand is at least
$1-d_X/m\ge0$. If $t=3$, then $m\ge2\cdot5\cdot7=70$ and the
three distinct indices $m/d_X>1$ are prime to three. Their
reciprocals sum to at most $1/2+1/4+1/5$, so the three summands
contribute at least $41/20>2+2/m$, contradicting the identity.
\end{proof}

\begin{lemma}[An isolated quartic prime]\label{fourm:heightone}
Let $q$ be isolated in an ordinary minimal table with $n=q^2N$, $q\nmid N$. If at least two labels have $q$-valuation one in their degree, there is a $q$-active label of degree $qN$.
\end{lemma}

\begin{proof}
Write $A_0,A_1$ for the active labels of degrees $N,qN$. These are all
the possibilities by Lemma~\ref{lem:complement}. Put $m=|A_0|$,
$\delta_i=\sum_{Z\in A_i}(1-\theta_{Z,q})$ and $b=q^2W_q e_0$.
Proposition~\ref{deep:cluster} makes $b$ integral,
$v_q(b_X)=v_q(d_X)$ for $v_q(d_X)<2$, and all height-one labels
$q$-principal. The vacuum Gauss formula gives, for a root of unity $\lambda$,
\[
      \lambda b_0=q^2-\delta_0/q^2-\delta_1 .
\]

For such a label $X$, every $Z\in A_0$ has $S_{XZ}\ne0$ by
principality, and $S_{XZ}/(qN)$ is integral by lcm divisibility.
Their Galois-mean contribution to the normalized row norm is at least
$m/q^2$ besides the vacuum; thus $m<q^2$. Since $b_0$ is a $q$-unit
and $\delta_1$ vanishes above $q$, the Gauss formula gives $m>0$.
All conjugates of the algebraic integer $1-\delta_0/q^2$ are strict
convex averages of $1$ and nontrivial roots of unity, hence have
modulus less than one. Therefore $\delta_0=q^2$. The mean of
$1-\theta$ for a nontrivial $q$-power root is at most $q/(q-1)$,
so $m\ge q(q-1)$. If $A_1$ were empty, $|b_0|=q^2-1$;
with the two height-one coordinates this costs at least
$(q^2-1)^2+2q^2>\|b\|^2=q^4$.

\end{proof}

\section{Three prime divisors}\label{short:three-prime}
Suppose that the dimension $M$ of a minimal obstruction has exactly
three prime divisors $p<q<r$. We first show why the two smaller primes
must both have exponent at least five. This reduces the entire range below
$5400^2$ to two local questions, one at three and one at two.

\begin{lemma}\label{short:quartic}
If $M=p^a q^b r^c$ is a minimal obstruction, then
$a,b\ge5$ and $c\ge4$.
\end{lemma}
\begin{proof}
Corollary~\ref{prop:reduced-quartic-two-exclusion} excludes exponents at most three.
Suppose a prime $h$ of exponent four has a larger companion $k$,
and call the remaining prime $t$. In a reduced table
$h=2$ is excluded by Corollary~\ref{prop:reduced-quartic-two-exclusion}.

First $h,k$ cannot occur together in a twist. Take a degree $ht^i$ in
$P_h$. Any $t$-free nonunit in $P_k$ has degree $h^j k^l$, $j=1,2$
(Lemma~\ref{three:complementary}). For $j=1$ the normalized crossing
has weight $h$ and residue $h$ at both $h,k$, contrary to
Theorem~\ref{three:word} and strictness. Thus the latter labels have
$h$-unit codegree and vanish on $h$-active columns. Were such a column
also $k$-active, orthogonality on $P_k$ modulo $t$ would leave only the
unit. This works in ordinary data, and also at odd $h,k$ in the
positive quotient.

In ordinary data $h$ cannot be isolated. Isolation makes the dimension
a square by Lemma~\ref{lem:complement}. The two complementary
degrees of $h$-height one in $P_h$ then give an active row of codegree
$h^2$ by Lemma~\ref{fourm:heightone}, forbidden by
Lemma~\ref{prop:quartic-prime-top-codegree}. If $h$ is the least
prime, separation from both larger companions already gives a
contradiction. We are left with $h,k$ odd, in either setting.

Work in the ordinary table or in the integral modular extension.
In the latter the remaining prime is $t=2$. Separation on the even
sector implies separation everywhere. Indeed, the space of joint $h,k$-active
labels is invariant under the dyadic Weyl which interchanges the
fermion permutation $F$ and monodromy sign $G$. A nonempty such space
supported on odd labels has no $G$-fixed vector but does contain the
$F$-fixed vector $v_X+v_{gX}$. Hence no character field ramifies at
both $h,k$ by Proposition~\ref{deep:cluster}.
In the extension we can lift odd-principal pairs to full principal
representatives by Lemma~\ref{lem:integral-minimal-extension}.

Choose a joint $t,h$-active column $J$. Such a column exists in the
ordinary case because $h$ cannot be isolated. In the extension the
two representatives of an even $h$-active pair have opposite twists. Choose
$A\in P_h$ of degree $hk^j$ and $B,C\in P_k$ of degrees
$h^v k^l,t^w k^m$, respectively, with all powers positive.
Orthogonality on $P_t$ modulo $k$ gives a $k$-free noninvertible
$V$ hitting $J$. In the extension $J$ is even, so the two
invertibles together contribute two. Its $h$-height is at most one,
since an $h$-unit codegree vanishes at $J$. It hits both $A,C$ by
principality, so coprime vanishing forces $d_V=t^i h$, $i>0$.
At $V,A$ the normalized entry is nonzero by $t$-principality,
has weight $h$, and vanishes modulo $h$ by $h$-principality.
At $V,B$ it has weight $h$ and positive residue $h$ at both $t,k$.
Proposition~\ref{deep:positive-inertia} forces the first entry
to ramify at $h$ and the second at $k>h$. Both belong to $K_V$,
contradicting separation.
\end{proof}

The inequality $p^5q^5r^4<5400^2$ forces $p=2,q=3$ since $2^55^57^4>5400^2$. It then gives $r=5$ or $7$.
At $r=7$ only $(a,b,c)=(5,5,4)$ fits. At $r=5$, divide by
$2^53^55^4$ to see that the integer multiplier is strictly below
$6$. Its possible values are $1,2,3,4,5$. All these candidates have $v_3(M)=5$
except $2^53^65^4$.

\subsection{What a quintic at three forces}
Write $M=2^a3^5r^c$, with $r=5$ or $7$. We show first that a minimal
obstruction of this form must be ordinary with $a,c$ even. This
handles the reduced and ordinary cases with the same calculation.

\begin{lemma}[The missing ternary phase]\label{short:ternary-phase}
In an ordinary or positive reduced minimal table of odd ternary
valuation, the ternary phases cannot all belong to
$\{1,\eta\}$ for one primitive cube root $\eta$.
\end{lemma}
\begin{proof}
Suppose they do. Write $U,W$ for the ternary unipotent and Weyl,
$P$ for the phase-one projection, $w$ for the complementary Weyl
image of the vacuum, and $s$ for the normalized table. In the
positive space use the symmetrized vacuum pair. Then $Pw=w$,
$Ww=sv_0$, and the ternary torus at $u$ acts on $w$ by the
Legendre symbol $(u/3)$. Bruhat gives
$PWU^tW^{-1}w=(-t/3)PWw$ for $t=1,2$ since
the outer upper unipotents disappear under $P$ and on $w$, and
$W$ carries the torus to its inverse. Applying the Fourier projector
of the missing nontrivial phase to $W^{-1}w$ now gives
$w=-G Psv_0$ for a quadratic Gauss sum $|G|^2=3$.

In the ordinary table the complementary-Weyl indicator
identity \eqref{eq:indicator-entry} therefore gives
$b_X=3\Delta^{-1}\sum_{Y\ {\rm inactive}}d_Y^2\chi_Y(X)$
as an algebraic integer when the indicator exponent is zero on the
ternary level and one on its complement. In the positive ring the
same expression is the sum of the indicators of $X,gX$ in the
extension. Indeed the odd terms cancel, and in the symmetrized
pair basis the complementary Weyl image has coordinates
$2(W_{3'}v_0)_Y$, since $W_{3'}$ exchanges the fermion permutation
and monodromy sign. In either case the expression is rational by
Galois stability. Orthogonality shows that the integral virtual
element $b=\sum_X b_X[X^*]$ has eigenvalues three on the inactive
rows and zero elsewhere. It is self-adjoint with unit coefficient one
and dimension three.

The symmetric integral multiplication matrix $T=N_b$ satisfies
$T^2=3T$. A diagonal entry $h=0,3$ has no neighbors, and at $h=1,2$
the off-diagonal entries are two signs, since $h(3-h)=2$. The unit
component is thus a cycle. Length at least five is impossible because
two vertices at distance two have a unique two-step path. In a
triangle the trace is a multiple of three, the vacuum diagonal is
one, and all diagonals are one or two. All three must equal one. The
matrix has rank one, forcing all three dimensions to be equal to
one by the positive dimension eigenvector. In a four-cycle let $A$
oppose the vacuum and let $Y,Z$ be its neighbors. Adjacent
diagonals sum to three. The dimension eigenvector equation at
$Y,Z$ gives $d_Y,d_Z\in\{d_A-1,d_A+1\}$. Coprime vanishing implies
$\chi_A(Y)=\chi_A(Z)=0$. But $b=1\pm Y\pm Z$ would then have
eigenvalue one, a contradiction in either table.
\end{proof}

\begin{lemma}[Ternary inversion]\label{short:ternary-inversion}
Suppose the obstruction has dimension $2^a3^5r^c$, $r>3$.
In its ordinary table or its integral modular extension there is
no noninvertible
$U\in P_2$ of degree $2^i3$.
\end{lemma}
\begin{proof}
Lemma~\ref{lem:integral-minimal-extension} gives strictness for
full noninvertibles and excludes prime-power degrees there.
Put $E=P_3$ in the ordinary case and
$E=P_3\cap\{B:\theta_{B,3'}=1\}$ in the extension. In the latter,
evaluation at $g$ makes $P_3$ even and the complementary-twist
condition excludes $g$. In both cases the unit is the only
invertible in $E$. The ternary block orthogonality and the
odd-valuation slice apply on $E$, since its projector commutes
with the ternary group.

Assume such a $U$ exists. Each odd-degree nonunit $B\in E$ has
degree $3r^j$, by the principal height bound at valuation five.
The normalized entry $z=S_{UB}/\lcm(d_U,d_B)$ has weight three,
strict house below three, and residues $1,0$ at all primes above
$2,3$, respectively. We claim that
$z=\lambda(\omega-\omega^{-1})$, with $\omega$ a primitive cube
root and $\lambda$ a dyadic root. The $3$- and $r$-primary root
configuration is unimodular \cite[Theorems~3 and~13]{BH}. A dyadic
basis obtained by choosing one of each pair of opposite roots makes
the full configuration disjoint signed copies of it. Thus a minimum
integral signed word for $z$ has length at most three. Erasing the
$3$-parts leaves a sum divisible by three in an unramified ring.
If zero it is a cancelling pair (the complementary orders exclude
a vanishing triple). Otherwise equality in the norm bound and
triangle inequality makes it three times a root. The word is
therefore $\lambda(\alpha-\beta)$ or
$\lambda(\alpha+\beta+\gamma)$ for ternary roots. Inverting the $r$-part and reducing at two gives
$\lambda_r^{2}=1$ in the residue field, since both reductions are
one. Reduction is faithful on odd-order roots, so $\lambda$ is dyadic.

Modulo two the ternary coefficient vector after subtracting one
is constant on each coset of $\mu_3$ by the cyclotomic relation.
For the pair case its parity support is of odd size at most three,
hence it is $\{1,\omega,\omega^{-1}\}$. For the triple case its
support is even of size at most four, hence empty. The word becomes
$\lambda(1+2\xi)$, $\xi\ne1$. If $\xi$ has order at least nine,
the character of $BB^*$ at $U$ has a full-weight positive
expression supported on
$T=\{1,\xi,\xi^{-1}\}$ and the corresponding Galois translate has support $T^4$. To expose each set, give each occupied $\mu_3$-coset
value one at its mark and $-1/2$ elsewhere,
zero on unmarked cosets, and tensor with normalized trace on the
other prime factors. This rational functional is at most one on
roots, with equality precisely on the marks. The adjoints of $B$
and its translate share a noninvertible, since otherwise rigidity makes
their tensor product simple, violating strictness of its double
braiding. In the extension $B$ is even and $gB\ne B$, so neither
adjoint contains $g$. The character of the common constituent at $U$ has balancing phases
which lie
in $T\cap T^4=\{1\}$ by positivity and exposure, again violating
strictness. Hence $\xi=\omega^{\pm1}$ and $1+2\xi=\pm(\omega-\omega^{-1})$. This proves the claim.

Let $\iota$ invert the ternary factor only and put
$H=\chi_U+\chi_{U^\iota}$. Square covariance gives $U$ and $U^\iota$ the same ternary phase.
The claim gives $H(B)=0$ and
$B^\iota\ne B$ on odd-degree nonunits in $E$. On its even-degree members $H(B)$
vanishes above two by principality.
If an odd-degree $J$ has ternary phase different from that of $U$,
block orthogonality gives $\sum_{B\in E}\chi_J(B)\overline{H(B)}=0$.
The vacuum contributes two. Every other nonzero term has dyadic
valuation greater than one, since lcm divisibility makes
$\chi_J(B)$ even when $d_B$ is even. Thus every odd-degree label has the phase of $U$, which equals one
by comparison with the vacuum.
The odd-degree members of $E$ number one modulo two, the nonunits
being paired by $\iota$. The slice identity on $E$ makes the mass
at primitive cubic phases twice the mass at phase one. Modulo four
these sides equal zero and two, respectively.
\end{proof}

\begin{proposition}
A minimal obstruction of dimension $2^a3^5r^c$, $r>3$, is ordinary
with $a,c$ even. There is a primitive cube root $\eta$ such that
the labels with ternary phases outside $\{1,\eta\}$ form a nonempty
set. Each has formal codegree $27$ or $243$.
\end{proposition}
\begin{proof}
Work in the ordinary table, or in the full extension.
An $r$-free noninvertible of $P_2$ has degree $2^i9$. Indeed prime powers are impossible, its degree squared divides the full
dimension, and Lemma~\ref{short:ternary-inversion} excludes
ternary height one. Each even dyadically active column in the
extension hits such a label by orthogonality on $P_2$ modulo $r$, since
the invertibles $1,g$ contribute two together. There is such a
column since the twists of the representatives of an even pair
are opposite. In the ordinary case the same argument uses the
single vacuum on every dyadically active column.

A $2$-free nonunit in the ordinary or positive $P_r$ has degree
$9r^j$. Indeed, if its ternary height were one, pair it with an
$r$-free nonunit of $P_3$, of degree $3\cdot2^k$ by
Lemma~\ref{three:complementary}. The crossing has weight three
and positive residue three at $3,r$, contradicting
Theorem~\ref{three:word} and strictness. These $P_r$ witnesses
exist by prime avoidance and lift to full principal representatives
in the extension, by Lemma~\ref{lem:integral-minimal-extension}.

Both kinds of witnesses now have ternary defect one. Since $r\nmid2$,
Lemma~\ref{lem:quintic-common-class} says that they all see only
$1$ and a single primitive cubic phase $\eta$. In the extension
one representative of every even pair is dyadically active, and
therefore hits a witness. Both representatives have the same ternary
phase. This contradicts Lemma~\ref{short:ternary-phase} in the
positive table.

In the ordinary case every $2$- or $r$-active column hits a witness
of the respective kind by prime avoidance. There is a ternary
phase outside $\{1,\eta\}$ by Lemma~\ref{short:ternary-phase}.
On each such fiber the complementary local groups act by scalars because their upper
unipotents do so and their conjugates generate them.
Lemma~\ref{lem:complement} implies that every label $Z$ there has
codegree $3^{5-2j}$, so $a,c$ are even. Codegree three is impossible.
Orthogonality to the vacuum reduced at each of $2,3,r$ requires three
distinct nonunit hits, since a degree omitting two primes would be
a prime power. Along with the vacuum these cost at least four in the
Galois-averaged row norm.
\end{proof}

Only the ordinary dimension $2^63^55^4=243n^2$, $n=200$, survives
this proposition. We exclude it by counting the exceptional ternary
labels in two ways.

\begin{proposition}
There is no ordinary minimal obstruction of dimension $2^63^55^4$.
\end{proposition}
\begin{proof}
There is a degree-$12$ label $U\in P_3$. Indeed a nonempty active
dyadic phase fiber contains a $3$-active column $J$. Otherwise
the ternary upper unipotent would be the identity on that invariant
fiber, as would its conjugates. They generate the local group, but
the torus of a nonsquare unit acts there by minus a permutation because
its Galois automorphism negates the square root of the dimension. The column
$J$ hits a $5$-free nonunit of $P_3$ with degree $3\cdot2^i$ by
prime avoidance. Here $i\ge2$ by Lemma~\ref{three:complementary},
and $i\ne3$ by dyadic defect-zero vanishing at $J$.
Every $5$-principal $3$-free nonunit has degree $4\cdot5$ or
$8\cdot5$ by the same lemma. Pair it with any degree-$12$ member
of $P_3$. The resulting normalized entry has weight four, strict
house below four and positive residue four at both $3,5$.
Theorem~\ref{three:word} gives an orbit of length four under tame
$5$-inertia, including in its signed alternative.

Suppose an exceptional row $Z$ has codegree $27$, so $d_Z=3n$.
All its hits lie in $P_3$ by Lemma~\ref{thm:principal-support}.
Apply the dimension trace identity \eqref{eq:r8-tracedim} modulo
two and five. It requires distinct nonunit trace-contributing
orbits of odd and of $5$-free degree, respectively. Each such
degree has ternary height one by the principal height bound and
exclusion of prime powers. For a representative $X$,
$\chi_X(Z)\equiv d_Z=0$ above three and
$\chi_Z(X)=(d_X/d_Z)\chi_X(Z)$ has the same congruence.
Its absolute field trace is therefore divisible by three. Put
$F=\mathbb Q(\chi_Z(X))$. By \eqref{eq:r8-tracenorm} the orbit
costs at least $9[K_X:F][K_Z:F]$. The budget $27$, including the
vacuum, forces both field indices to equal one, both traces to be
$\pm3$, and uniqueness of the two required orbits of the indicated
types. The $5$-free one has degree $12$ or $24$ by the complementary
bound. Degree $12$ is ramified at five, whereas $5\nmid f_Z$ makes
$K_Z$ unramified there by Lemma~\ref{cor:field-conductor}. Degree $24$ would require $1\pm3\cdot24=0\pmod5$.
Thus every exceptional label has codegree $243$ and degree $n$.

Let $L$ count labels at ternary phases of order at least nine
and $H$ those at the primitive cubic phase other than $\eta$.
Square Galois covariance makes the higher phases equidistributed
over each orbit of square units. Such an orbit is a union of
complete $\mu_3$-cosets, so these labels contribute zero to the
ternary Gauss sum $T=\sum_Xd_X^2\theta_{X,3}$.
In units of $n^2$, the masses at the trivial and the two primitive
cubic phases are $A=(243-L)/3$, $2A-H$, and $H$, respectively,
by the odd-valuation slice. Consequently
$|T|^2=3n^4(A-H)^2$. The indicator of the regular object at an
exponent selecting the ternary twist is $|T|^2/(243n^2)$
by summing \cite[Theorem~7.5]{NSfs} over the simples with their dimension multiplicities. This is a rational algebraic integer.
Thus $9\mid A-H$, or $27\mid L+3H$. Since exceptional labels
exist, their number $N=L+H$ is at least nine.

Every exceptional label hits $U$ by $3$-principality, and its
field is unramified at five. The corresponding unnormalized
entries on the orbit of $U$ are divisible by $\lcm(12,n)=3n$.
Averaging over that inertia orbit and all embeddings as in
Lemma~\ref{deep:field-energy} gives $4N(3n)^2<243n^2$, impossible
for $N\ge9$.
\end{proof}

\subsection{The remaining dyadic quintic}

It remains to exclude a minimal obstruction of dimension
$2^5 3^6 5^4$. Write $n=3^3 5^2$. We shall use the same
complementary witnesses in both tables: the $5$-free nonunits of
$P_3$ and the $3$-free nonunits of $P_5$. By
Lemma~\ref{three:complementary} they have degrees $4\cdot3^h$ and
$4\cdot5^l$, respectively. Indeed their two-part is at least three
and its square divides the table dimension, whose two-part is $32$
or $16$. Orthogonality with the vacuum on $P_s$, reduced at the
omitted prime, makes every $s$-active column hit one of the
corresponding witnesses ($s=3,5$).

\begin{lemma}[Dyadic phases at odd valuation]\label{seq:two-phase}
In an integral modular table with $v_2(D)$ odd, a row of dyadic
defect one sees only the phases $1$ and $i$ for one primitive
fourth root $i$. Its full two-block has one label of its degree
at each phase. If all dyadic twists of the table are $1,i$, the
category has a nontrivial invertible.
\end{lemma}
\begin{proof}
Let $E$ project onto a full two-block $B$ and $Q=sEs^{-1}$.
Both commute with the dyadic group. Work modulo the dyadic part $2^a$
of $\ord(T/\gamma)$. This is divisible by eight since the level field
contains $1/s_{00}=\sqrt D$ and its quadratic discriminant has
two-part eight. Put $U(1)=\lambda T_2$, with $\lambda=\gamma^{-1}$,
$\lambda^8=1$ by integrality \cite[Proposition~6.7(ii)]{DLNCongruence}.
For odd $t$ set $u=1+4t$ and $M(t)=\sum_{j\in B}|s_{0j}|^2
\theta_{j,2}^{-2t}$. The dyadic torus $H_u$ negates the vacuum:
$u\equiv5\pmod8$ negates $\sqrt D$ and the complementary Galois
action is trivial. Since $L(b)=sU(-b)s^{-1}$ we have
$(QL(2t))_{00}=\lambda^{-2t}M(t)$. The vacuum entry after $Q$ in
$U(2)L(2t)=L(2t/u)H_uU(2/u)$ therefore reads
\[
       \lambda^{2-2t}M(t)=-\lambda^{(2-2t)/u}M(t/u).
\]
The exponents differ by $8t(1-t)/u\equiv0\pmod8$.
The map $t\mapsto t/(1+4t)$ permutes odd residues modulo $2^a$
with inverse $h\mapsto h/(1-4h)$. Thus the mean of $M(t)$ is
zero. In this mean phases of order one or two contribute their
squared mass, phases of order four its negative, and higher
orders zero.

Every block contains a phase-one label: $W_{2'}^{-1}$ preserves
dyadic twists, and $s^{-1}\CC^E=W_{2'}^{-1}\CC^E$ has
nonzero vacuum coordinate. In a defect-one block there are at most
two inertia orbits by Lemma~\ref{lem:one-defect-block};
they must have phases $1$ and $i$, respectively, since inertia
preserves the order and fixes each fourth root under square action.
The given row is unramified, as its orbit would otherwise exhaust
the block. On this block $U(4),L(4)$ are scalar. The identity
$U(1)L(4)=L(4/5)H_5U(1/5)$ makes $H_5$ scalar, since
$U(1/5)$ differs from $U(1)$ by a scalar. The units $1\bmod4$
are generated by five, so the other orbit has at most two labels.
Equality of their masses and integrality of degrees force two
singletons of equal degree. The support assertion follows again
from $s^{-1}\CC^E=W_{2'}^{-1}\CC^E$.

If the whole table has just the indicated phases, set
$Z=T_2^2$, $F=sZs^{-1}$ and $P=(1+Z)/2$. As above $H_5=-I$,
and $U(2)L(2)=L(2/5)H_5U(2/5)$ yields $ZF=-FZ$, hence
$PFP=0$. Write $T_2=uI+vZ$ with $u=(1+i)/2$, $v=(1-i)/2$.
Up to a scalar phase the inverse dyadic Weyl is
$T_2(uI+vF)T_2$, so $PW_2P$ on the phase-one space is scalar
of modulus $1/\sqrt2$. For $w=W_{2'}e_0$ we have $Pw=w$,
$W_2w=se_0$. Hence $|w_j|^2=2d_j^2/D$ on phase one and zero
elsewhere. At an indicator exponent zero on the dyadic level and
one on the odd level, \eqref{eq:indicator-entry} gives an integral
virtual element with eigenvalue two on phase one and zero elsewhere.
Indeed the complementary upper--lower--upper identity gives the
squared entries $|w_j|^2$, and the resulting indicators are rational
by Galois stability. Orthogonality gives unit coefficient one, squared
norm two and dimension two. The remaining coefficient is therefore
a positive invertible.
\end{proof}

\begin{proposition}\label{short:ordinary-dyadic-quintic}
For $n=3^3 5^2=675$ there is no ordinary minimal table of dimension $32n^2$.
\end{proposition}
\begin{proof}
All complementary witnesses have dyadic defect one. Opposite witnesses
hit both equal-degree representatives in each other's two-block by
odd-prime principality. The preceding lemma thus gives one common
visible pair of phases $1,i$ and at least one label outside it.
An outside fiber is inactive at both odd primes. Their local groups
preserve the fiber and their upper unipotents and conjugates act by
scalars there. Thus each label $Z$ in it has pure dyadic codegree
by Lemma~\ref{lem:complement}: its degree is $2^jn$ and its codegree
is $2^{5-2j}$, $0\le j\le2$.

The row is supported on $P_2$ (Lemma~\ref{thm:principal-support}).
Reducing its dimension-trace identity at three and five forces
distinct nonunit orbits with nonzero trace of degrees $2\cdot3^h$
and $2\cdot5^l$: a prime-power degree is impossible and principal
two-height is at most one. Their traces and the vacuum already exceed
codegree two. At codegree eight one has $j=1$. For a representative
$X$ of either orbit, $\chi_X(Z)$ vanishes above two by principality,
as does $\chi_Z(X)=(d_X/d_Z)\chi_X(Z)$ because the ratio is a
two-adic unit. Its absolute trace is even; the two contributions
cost at least $4+4$ in the trace norm, besides the vacuum
(Lemma~\ref{deep:field-energy}). Hence every outside label has
degree $n$ and codegree $32$.

Let $N$ count these labels. If the largest dyadic twist order is
$2^s\ge8$, let $L$ count labels at that order. The Gauss sum for
the exponent $2^{s-1}$ on the dyadic level and zero elsewhere is
$n^2(32-2L)$. Its squared modulus divided by $32n^2$ is the
regular-object indicator, a rational algebraic integer. Thus
$4\mid L$ and $N\ge4$. If the level divides four, phases
$1,-1$ together have mass $16n^2$, as do the phases $i,-i$, by
the first paragraph of the lemma. For $u,v$ the respective counts
at $-1,-i$, the dyadic Gauss sum is
$n^2((16-2u)+i(16-2v))$. Integrality gives $8\mid u^2+v^2$.
Squares modulo eight show $4\mid u+v$, so again $N\ge4$.

Take a $5$-free nonunit $X\in P_2$ of degree $2\cdot3^h$ by
Proposition~\ref{prop:principal-support}. At the opposite witness
$4\cdot5^l$ the normalized entry has weight two and residue two
at five. Theorem~\ref{three:word} makes it a unique unsigned
five-primary word of length two. Strictness requires a nonidentity
term; the tame stabilizer acts freely on such terms, so the
five-inertia orbit of $X$ has length $e\ge2$. All outside labels
are unramified at five and hit $X$ by two-principality and odd
degree. Their normalized entries agree along this orbit. The field
energy bound (Lemma~\ref{deep:field-energy}) gives
$eN(2n)^2<32n^2$, which contradicts $e\ge2$, $N\ge4$.
\end{proof}

For the reduced case put $\Delta=16n^2$, and let $\M$ be the
integral minimal modular extension of dimension $D=4\Delta$.
Its even sector is the given category, with fermion $g$. Write $F$
for tensoring by $g$, $G$ for its monodromy sign, and $W,U=\lambda
T_2$ for the dyadic Weyl and upper unipotent. We use the three spaces
\[
 A=(F=G=1),\qquad B=(F=-1,G=1),\qquad C=(F=1,G=-1).
\]
Thus $W$ preserves $A$ and interchanges $B,C$; the odd-primary
Weyl $W_o$ commutes with $F,G$. An even fermion orbit contributes
one basis vector to each of $A,B$, and every odd orbit contributes
a vector to $C$.

\begin{lemma}[An isolated squared dyadic phase]\label{short:deep-pair}
Call a positive pair deep if its squared dyadic twist differs from one.
If all pairs in a deep squared-twist fiber are inactive at every odd
prime, they have pure dyadic positive codegree and belong to the
positive principal two-block. If $Z$ is the only deep pair in the
table and its degree is odd, then under this hypothesis
$v_2(f_Z^+)\le2$.
\end{lemma}
\begin{proof}
On this fiber in $A$, each odd local group acts by scalars: its
upper unipotent does so, and the conjugate unipotents generate the
group. Its Weyl preserves each symmetrized pair line. Commutation
with $T_2$ separates the two full label lines, since their dyadic
twists are opposite. The line criterion of Lemma~\ref{lem:complement}
thus gives pure dyadic codegree.
The dyadic Weyl $K=W|_A$ preserves positive two-blocks: their full
lifts are blocks by the squaring argument of
Lemma~\ref{lem:fermion-quotients}, and the local Weyl preserves
these blocks. As the odd Weyls preserve the pair line, its $K$-column
agrees up to phase with its normalized reduced $S$-column. The
vacuum entry is nonzero, proving principality.

If $Z$ is unique, a power of the squared dyadic twist is
$R=I-2P_Z$ on $A$, where $P_Z$ projects onto that pair.
The commutator $RKR^{-1}K^{-1}$ has finite order; scalar lifts of
the local unipotent cancel. Its trace is
$\dim A-4+4|K_{ZZ}|^2$, so $4|\chi_Z(Z)|^2/f_Z^+$ is integral.
For odd $d_Z$ the factor $|\chi_Z(Z)|^2$ is a dyadic unit by
principality, which gives the bound.
\end{proof}

\begin{lemma}[When even dyadic twists are signs]\label{fam:Q:5.9}
In an integral modular extension of dimension $D=4\Delta$ as above,
if all even dyadic twists are signs, then $v_2(\Delta)$ is even.
Under the same sign hypothesis, if $D$ is a square and $8\mid\Delta$, no positive principal
two-block character has dyadic defect two.
\end{lemma}
\begin{proof}
By Lemma~\ref{lem:signed-unramified-sign-twists} dyadic inertia
fixes even labels. The torus $W^2$ acts on $A,B$ by the Galois
sign $\epsilon$ of $\sqrt D$, and on $C=WB$ by the same sign.
Put $J=W|_A$ and $P=WU:A\to C$; this is a bijection as $U$
interchanges $A,B$. On $A,B$ one has $U^2=\lambda^2$. Applying
$WUW=U^{-1}WU^{-1}$ to $Ux$, $x\in A$, gives
$UPx=\epsilon\lambda^{-2}PJx$. Since $J^2=\epsilon$, the odd
dyadic phases belong to a pair $\delta,-\delta$. In the full
Gauss sum the even pairs cancel. Thus
$p^+/\delta=(\xi/\delta)\sqrt D$ lies in an odd cyclotomic field
$K$, with $\xi$ the central charge. Squaring shows that the square
of $\xi/\delta$ lies in $K$, so its odd part belongs to $K$ and
its dyadic part has order at most four. Hence $\sqrt D\in K(i)$.
If $v_2(D)$ were odd, the quadratic discriminant would be divisible
by eight, which is impossible in a field of conductor dividing
four times an odd number. This proves the parity assertion.

Now suppose $D$ is a square. Choose $k$ zero on the dyadic level
and two on the odd level. For even $V$, the two inputs in the
fusion formula for $\nu_k(V)$ have the same fermion sector, hence
their dyadic twist ratio is a sign; thus $\nu_k(V)=\nu_2(V)$.
For $b=\sum_V\nu_k(V)[V^*]$, equation~\eqref{eq:indicator-entry}
gives $\chi_x(b)=f_x|(sT^ks^{-1})_{0x}|^2$. Odd-local Bruhat at
the unit parameter $-2$ identifies this absolute entry with one
of the odd Weyl vacuum column, up to an odd torus permutation.
That column is rational by Lemma~\ref{lem:local-vacuum-column},
and is supported on even labels with dyadic twist one. The torus
preserves both conditions. Thus $gb=b$ and at most one of
$\chi_x(b),\chi_{gx}(b)$ is nonzero on an even pair.
In the positive quotient
$V_+=\sum_{[V]}\nu_2(V)[V^*]$ has eigenvalues
\[
 w_x=\tfrac14(\chi_x(b)+\chi_{gx}(b)).
\]
Indeed averaging keeps the even terms and $gb=b$ makes the two
coefficients in each pair equal. Each $w_x$ is a rational square
as $f_x=D/d_x^2$ is a square in $\QQ$. It is an algebraic integer
as an eigenvalue of an integral matrix, hence an integer square.

Modulo two, $V_+^2$ equals the Casimir element
$\sum_{[V]}[V][V^*]$: nonselfdual pairs cancel, and a self-dual
pair has a self-dual representative since $gV$ and $V^*$ have
opposite twists. Consequently $4\mid w_x$ if the positive codegree
$f_x^+$ is even. Suppose $x$ were positive principal of defect two.
Because its field is unramified at two,
$h(V)=(\chi_x(V)-d_V)/2$ is locally integral. Character
orthogonality gives
\[
 \sum_{[V]}h(V)h(V^*)=(f_x^++\Delta)/4\equiv1\pmod2;
\]
here $x$ differs from the vacuum since $8\mid\Delta$.
Nonselfdual pairs cancel, so by Frobenius the sum
$\sum_{[V]}\nu_2(V)h(V^*)$ is nonzero modulo two. It equals
$(w_x-w_0)/2$, which is even.
\end{proof}

\begin{lemma}[The odd sector has two phases]\label{seq:fermion-one-spin}
If all even dyadic twists are signs in the minimal extension,
the odd dyadic twists cannot have a common value.
\end{lemma}
\begin{proof}
Suppose $U$ is constant on $C$. The identity for $UP$ in the
preceding proof then makes $J$ scalar on $A$. Write $w=W_oe_0$.
It is supported on even labels of dyadic twist one.
Projecting $se_0=Ww$ onto $A$ shows that
$|w_j|^2=4d_j^2/D$ on the set $E$ of those labels: $W$ is
scalar on $A$ and symmetrization adds the two equal vacuum entries.
The complementary indicator calculation in Lemma~\ref{seq:two-phase}
gives a self-adjoint integral virtual element $b$ with eigenvalue
four on $E$ and zero elsewhere. On the even basis its coefficients
are those of $1+g$. Indeed one label per even pair lies in $E$,
and the two have equal characters on the even sector; the sum
over the even sector gives the character projector for $g=1$.
Hence $v=b-1-g$ is supported on odd labels, $gv=v$, $\dim v=2$
and $v^2=2(1+g)$ by character evaluation. Its norm squared is two.
A free fermion pair in its support would give an odd invertible;
hence $v=Y-Z$ for two fixed odd simples, which are self-dual.
Then $vY$ has squared norm four, unit and $g$ coefficients one,
and remaining part $\pm(H+gH)$ for an even noninvertible $H$.
By reciprocity $vH$ hits $Y$, has squared norm two and is
orthogonal to $v$. Positive dimension gives $vH=Y+Z$.
Together with $\dim v=2$, this yields
$d_Y=d_H+1$, $d_Z=d_H-1$. Both are coprime to $d_H$, so
$\chi_H(Y)=\chi_H(Z)=0$ by Lemma~\ref{lem:integral-minimal-extension},
contrary to $\chi_H(v)^2=4$.
\end{proof}

\begin{proposition}\label{short:fermionic-dyadic-quintic}
For $n=3^3 5^2=675$ there is no slightly degenerate minimal obstruction of dimension $32n^2$.
\end{proposition}
\begin{proof}
We first prove that all even dyadic twists are signs. The
complementary witnesses have odd positive codegree, so their
positive two-blocks are singletons. Block orthogonality with the
vacuum makes their rows vanish at all deep pairs. Each deep fiber
is therefore inactive at three and five by the witness property at
the start of this subsection. Its codegrees are pure dyadic by
Lemma~\ref{short:deep-pair}. Codegree one exhausts the dimension
and Lemma~\ref{prop:quartic-prime-top-codegree} excludes codegree
four; hence a deep pair has codegree $16$ and odd degree $n$.

If any such pair exists, orthogonality on positive $P_2$ with the
vacuum modulo five supplies a $5$-free nonunit $X$ of degree
$2\cdot3^h$. As in the ordinary proof, comparison with the
$4\cdot5^l$ witness puts $X$ in a five-inertia orbit of size
$e\ge2$. The $N$ deep pairs have odd degree and are
unramified at five, so $eN(2n)^2<16n^2$. The deep pair is unique,
contrary to Lemma~\ref{short:deep-pair} at codegree $16$. Thus
all even dyadic twists are signs, and Lemma~\ref{fam:Q:5.9}
and the principal valuation bound make all positive principal degrees odd.

Choose full principal lifts of the complementary $P_3,P_5$
witnesses as in Lemma~\ref{short:quartic}. Their positive
two-blocks are singletons. For each such $X$ its full two-block
consists of $X,gX,T,gT$ with $T$ odd, $d_T=d_X$. Indeed the
positive and signed orders have identical special fibers at two
(Lemma~\ref{lem:fermion-quotients}), hence the corresponding
signed block has one character. Agreement of residues on even
labels extends to odd labels by squaring. Comparing the $1,g$
columns on this full block equates odd and even squared masses.
A fixed odd label would require $d_T^2=2d_X^2$, so it is a free pair
of the asserted degree. Any odd hit of $X$ has the phase of this
mate: $W^{-1}$ preserves the block and $W_o^{-1}$ preserves sector
and dyadic twist. Every opposite witness hits the mate by odd-prime
principality, so all these odd hits have a common dyadic phase
$\delta$.

For $s=3,5$, a full principal $s$-label is even by evaluation
at $g$; $g$ itself is excluded by evaluation on an opposite odd
mate of $s$-unit degree. Orthogonality with the vacuum on $P_s$
modulo the omitted prime therefore makes every $s$-active odd
column hit a complementary principal witness. An odd dyadic
fiber distinct from $\delta$ consequently is inactive at both
odd primes, and each of its labels $Y$ has pure dyadic full
codegree by Lemma~\ref{lem:complement}. Such a $Y$ exists by
Lemma~\ref{seq:fermion-one-spin}.

Choose an odd mate $T$ of a $5$-free witness. Its phase is
$\delta$, so its signed character differs from that of $Y$.
Orthogonality on even pair representatives, reduced modulo
five, gives a $5$-free nonunit pair $V$ hit by $Y$: the vacuum
contributes one, and $5\mid\chi_T(V)$ whenever $5\mid d_V$ by
lcm divisibility. The pure dyadic codegree of $Y$ puts $V$ in the
full principal two-block by Lemma~\ref{thm:principal-support}.
Its positive pair is principal and hence of odd degree. That
degree can only be a power of three, contrary to coprime vanishing
and prime avoidance in the positive table.
\end{proof}

\section{Four prime divisors}\label{short:four}
Let $M<5400^2$ be a divisor-minimal obstruction with four prime divisors.
Its exponents are at least two. We first show that two of them are exactly
two, and that the obstruction is ordinary.

\paragraph{Locating the two squares.}
A cube $p^3$ and a square $t^2$ cannot both occur as exact prime powers.
Indeed, there is a label active at both $p,t$. In the ordinary case
separation would force a $t$-active label to have $p$-unit codegree,
impossible at odd $v_p(M)$. In the reduced case $t$ is odd since
$v_2(M)=2$ would leave a prime of valuation one in the table dimension.
If $p$ is odd, the nonempty $t$-active sector is invariant under the
$p$-local group and contains a $p$-active label. In fact a wholly
$p$-inactive invariant subspace would be fixed by both elementary
unipotents, whereas a nonsquare local torus acts there by minus a
permutation since $v_p(M/2)$ is odd. If $p=2$, work in the integral
minimal extension and choose between an even $t$-active label and its
fermion translate to obtain dyadic activity.
On the principal $p$-block the degrees are $p$-free
(\ref{cor:principal-valuation-bound}, or \ref{lem:principal-fermion-two}
in the extension at two). Those hit by the joint label are also
$t$-free by defect-zero vanishing. The noninvertible hits therefore
contain both remaining primes: in the minimal tables a nonunit has no
prime-power degree, and in the full dyadic block its noninvertible supports
intersect pairwise with empty total intersection
(Proposition~\ref{lem:principal-fermion-two}). Block orthogonality against the vacuum,
reduced at an odd one of those primes, leaves $1=0$, or $2=0$ in the
extension since the joint label is even.

If there were no square, the product of the four primes would be at most
$\sqrt[3]{5400^2}<308$, hence would be $210$. The dimension is then
$210^3$ times $1,2$ or $3$, so $5$ and $7$ both have valuation three.
On the principal seven-block take a nonunit of odd degree by prime
avoidance at a seven-active column. Its degree is seven-free and is
therefore supported on exactly $3,5$. Call this label $X$. On the
principal five-block all degrees are five-free. A nonunit $Y$ of
three-free degree would be coprime to $d_X$, hence miss $X$.
This contradicts prime avoidance on that block at $X$ with omitted
prime three: its odd-valuation alternative applies since $3\nmid5-1$.
With exactly one square, the other three exponents would now be at least
four, giving $M\ge(2^23^25^2\cdot7)^2>5400^2$. This proves the claim
about the two squares.

Write $M=p^Aq^B u^2v^2$ with distinct primes and $A,B\ge2$.
In the reduced case a squared two is again impossible, so take $p=2$
and $q,u,v$ odd. A $u$-active column hits a nonunit odd degree:
reduce orthogonality to the vacuum at two. Such a hit is $u$-free,
hence has degree $q^iv$, $i>0$. Similarly $q^{i'}u$ occurs. In the
integral minimal extension every principal two-block noninvertible of
$q$-free degree hits the lifts of both these odd degrees by principality,
and therefore contains both $u,v$ in its support. Choose an even
$u$-active label and a fermion translate making it two-active. On the
principal two-block its orthogonality to the vacuum, reduced modulo $q$,
retains only the two invertibles: a $q$-free noninvertible has
$u$-unit codegree. This is impossible. We henceforth work in the
ordinary case.

\paragraph{Consequences of isolation.}
At a $u$-active column the hits are $u$-free. Reducing vacuum
orthogonality modulo $p$ and $q$ supplies degrees $q^iv$ and $p^lv$,
respectively; similarly $q^{i'}u$ and $p^{l'}u$ occur. A $p$-free
nonunit of the principal $q$-block must hit the displayed degrees
$p^lv,p^{l'}u$, and is therefore divisible by $uv$. A $q$-active
column hits such a nonunit by prime avoidance, so cannot be active at
either $u$ or $v$. Interchange $p,q$. Nor can both $u,v$ divide one
twist order: every nonunit hit on that column would be divisible by
$pq$, contradicting vacuum orthogonality. Thus $u,v$ are isolated.
Each $t$-active label, $t=u,v$, has codegree $t^2$ by
Lemma~\ref{lem:complement}. In particular $M=n^2$ is a square, its
degree is $n/t$, and $u,v\ge5$ by Lemma~\ref{lem:isolated-dyadic-square}.
For $n=p^h q^k uv$, the $t$-free axis degrees supplied above include
one from each of
\[
 \{p^i w:1\le i\le h\},\qquad \{q^l w:1\le l\le k\},
 \qquad w=uv/t .
\]

We recall the arithmetic we will use from the local columns
(\ref{second:two-columns}, \ref{lem:singleton-energy}). For an
isolated prime $t\parallel n$, let $[b]_t$ denote the centered residue
modulo $t$, $j_t$ its deficit and $\epsilon=\pm1$ the common sign. Then
\begin{equation}\label{short:vacuumcheck}
 (n/t)j_t\equiv-\epsilon\pmod t,\quad 1\le j_t<t,\quad
 j_t\mid t-1\ \text{or}\ j_t\mid t+1,\quad
 \sum_t (n/t)j_t\le n+1 .
\end{equation}
At odd $n$ all $j_t$ are even. The residual squared norm at $t$
is at most $B_t=j_t(t+1-j_t)$; the mixed scalar product for two
singletons is $-j_uj_v$. Every proper $t$-free degree in the cases
below has column coordinate
\[
                     A_t(d)=[d/(n/t)]_t=[-\epsilon j_t d]_t .
\]
Indeed its coordinate has that residue, and there are at least two
other distinct proper $t$-free labels (Lemma~\ref{lem:isolated-dyadic-square}).
Thus its square is strictly below $(t+1)^2/4$, which forces centering.

\paragraph{Three or more singletons.}
Write $n=Huvw$, $u<v<w$ singleton primes, with $H$ a power of the
remaining prime. Applying the isolation argument to each pair
separates all four primes. For a pair $t,s\in\{u,v,w\}$ with third
member $z$, the only possible nonunit degree omitting $t,s$ is $Hz$.
Indeed $ts$ occurs by reducing vacuum orthogonality on a $z$-active
column modulo the prime of $H$. Any nonunit omitting $t,s$ has degree
supported on the other two primes; if its height at the prime of $H$
were deficient, full isolation would put it in that principal block
by Proposition~\ref{deep:cluster}. It would have to hit the coprime
degree $ts$. The mixed scalar product therefore gives the divisibility
\begin{equation}\label{short:pairdiv}
          |[s^{-1}]_t[t^{-1}]_s|\ \mid\ j_tj_s .
\end{equation}
For clarity the entire arithmetic check is given below. With a
nonsingleton power $H$, the inequalities $u\ge5$ and $Huvw<5400$
give $H=4,8,9$; if all four are singletons their only possible product
is $5\cdot7\cdot11\cdot13$, represented by $H=5$.
In the middle column first use $w<5400/(Huv)$. To obtain the last
column insert $j_t$ equal to the least positive residue of
$-\epsilon(Huvw/t)^{-1}$ in \eqref{short:vacuumcheck} for the two signs.
\[
\begin{array}{cc|l|l}
H&(u,v)&w&\text{surviving }(w;\epsilon;j_u,j_v,j_w)\\\hline
4&(5,7)&11,13,17,19,23,29,31,37&(29;+;2,1,6),\ (31;-;2,2,2)\\
4&(5,11)&13,17,19,23&(13;+;2,3,1),\ (17;+;3,1,1)\\
4&(5,13)&17,19&\varnothing\\
4&(7,11)&13,17&(13;-;3,1,3)\\
8&(5,7)&11,13,17,19&\varnothing\\
9&(5,7)&11,13,17&\varnothing\\
5&(7,11)&13&\varnothing
\end{array}
\]
In the displayed order, the five survivors violate \eqref{short:pairdiv}
as $6\nmid2$, $6\nmid4$, $10\nmid2$, $2\nmid3$, $6\nmid3$.
Use the pair $(u,v)$ except in the third test, where it is $(u,w)$.

\paragraph{Exactly two singletons.}
Both higher primes must be $2,3$: an odd root costs at least
$3^25^2\cdot7\cdot11$, and an even one omitting three costs at least
$2^25^2\cdot7\cdot11$. Write $n=2^h3^k uv$ with $5\le u<v$.
From $n<5400$ the possible triples $(h,k;uv)$ are exactly
\[
\begin{array}{c|l}
(2,2)&5\cdot\{7,11,13,17,19,23,29\},\quad7\cdot\{11,13,17,19\},\quad11\cdot13\\
(3,2)&5\cdot\{7,11,13\}\\
(2,3),(4,2)&5\cdot7 .
\end{array}
\]
We give a two-column test small enough to evaluate directly. A positive
multiplicity at degree $d$ requires
\begin{equation}\label{short:energycheck}
     j_t k_t(t-k_t)<t^2,\qquad k_t=|[d]_t|.
\end{equation}
If the character field is unramified at $t$, the stronger tests
\begin{equation}\label{short:unramcheck}
  A_t(d)=-\epsilon j_t[d]_t,\qquad (t-1)j_t[d]_t^2<t^2
\end{equation}
hold for a proper $t$-free label. Failure of \eqref{short:unramcheck}
therefore costs at least two labels of this degree, since a singleton
Galois orbit has rational field.

At each endpoint $t$ reserve one axis of each kind. For $b=2,3$ take
the least cost $A_t(d)^2$ among $d=b^i(uv/t)$ satisfying
\eqref{short:energycheck}, doubling the cost for each candidate failing
\eqref{short:unramcheck} before taking the minimum. Subtract these
two minima from $B_t$ to obtain $R_t$; use $-\infty$ if an axis is
empty. A common degree $x=2^i3^l$ must have an unramified endpoint
by Proposition~\ref{deep:cluster}. At that endpoint $t$ one also has
$t<2\min(2^i,3^l)$. Indeed cross $x$ with one of the two compulsory
$t$-free axes of gcd $g\le\min(2^i,3^l)$. Both labels are
$t$-principal. Their lcm-normalized entry $z$ has strict house $<g$,
lies in the unramified field of the first label, and is congruent to $g$
at every prime over $t$ by principality and Galois covariance.
If $t\ge2g$, the algebraic integer $(z-g)/t$ has house $<1$, giving
$z=g$, a contradiction.

The final column lists the possible signed products $A_u(x)A_v(x)$.
To check it, try $2\le i\le h$, $1\le l\le k$ (the preceding
inequality excludes $i=1$). Discard $x$ failing \eqref{short:energycheck}
at either endpoint. Count the endpoints satisfying both
\eqref{short:unramcheck} and $t<2\min(2^i,3^l)$. If the count is
zero discard $x$; if the count is one its multiplicity is at least
two. With this lower multiplicity $m$ (one when both qualify), require
$m A_t(x)^2\le R_t$ at both ends.
\[
\begin{array}{r|r|rr|rr|l}
 n&\epsilon&j_u&j_v&R_u&R_v&A_u(x)A_v(x)\\\hline
1260&+&2&4&4&0&\varnothing\\
1980&-&1&3&3&16&-2\\
2340&-&2&6&4&-\infty&\varnothing\\
3420&+&1&2&3&-16&\varnothing\\
4140&+&3&6&5&22&\varnothing\\
5220&-&4&5&4&-\infty&\varnothing\\
4284&+&2&6&4&-\infty&\varnothing\\
4788&-&3&4&-1&12&\varnothing\\
2520&+&1&2&3&8&1,2\\
3960&+&2&4&4&0&\varnothing\\
4680&-&1&3&3&-12&\varnothing\\
3780&-&1&1&3&5&1\\
5040&+&3&1&5&5&\varnothing
\end{array}
\]
Only these thirteen rows pass \eqref{short:vacuumcheck}. For example
at $n=1980$ the two axes cost at least $1,1$ at five and $9,2$ at
eleven. Of the two common degrees $12,36$ only $36$ fits, giving
$A_5(36)=1$ and $A_{11}(36)=-2$. Thus the mixed scalar product is
even although it must equal $-3$. In every other row it cannot be
negative. This excludes all four-prime obstructions.

\section{Five prime divisors}\label{sec:short-five}
Write $P$ for the product of the five primes of a minimal source $M<5400^2$.
Since $P<5400$, the primes $2,3,5$ all divide $P$ (even the least product
omitting $5$ is $2\cdot3\cdot7\cdot11\cdot13=6006$). Thus $P=30uv$ with
$7\le u<v$ and $uv<180$. The possibilities are
\[
              P=2310,2730,3570,3990,4290,4830 .
\]
As $M$ is square-full, $M/P^2<(5400/2310)^2<6$. We treat a prime
multiplier first, then the two square cases $M=P^2$ and $M=4P^2$.
The latter requires $P<2700$, hence $P=2310$.

\paragraph{A single cube.}
Write $M=p^3(u v a b)^2$. In an ordinary table the cubed prime $p$
mixes with each of $u,v$: otherwise a label active at that other prime
would have $p$-unit codegree by Lemma~\ref{lem:complement}, impossible
at an odd valuation of $M$. On a column active at $p,u$, avoidance on
$P_p$ gives a $b$-free nonunit hit (Proposition~\ref{prop:principal-support}).
Its degree omits $p$ by the principal valuation bound, and omits $u$
by defect-zero vanishing. It therefore has support exactly $\{v,a\}$
since prime-power degrees are impossible. On a column active at $p,v$,
interchange $u,v$ and $a,b$ to obtain a degree supported on $\{u,b\}$.
These two labels of $P_p$ have coprime degrees, so their crossing
vanishes; principality and $p$-freeness make it nonzero.

In a reduced table the square-fullness of $\Delta=M/2$ forces $p=2$.
In the integral modular extension the full principal two-block has odd
degrees, and its noninvertible supports intersect pairwise
(Proposition~\ref{lem:principal-fermion-two}). At each of $u,v$ take
an even active label and choose its fermion mate to make it two-active.
The same two avoidances hold: in vacuum orthogonality on the full
two-block the two invertibles contribute $2$, a unit at $a,b$.
The two hits have supports contained in the disjoint sets $\{v,a\}$
and $\{u,b\}$, contradicting pairwise intersection.

\paragraph{The two ordinary square cases.}
Put $M=n^2$. Call a prime a singleton when it occurs once in $n$.
If a degree has support $\{x,y\}$ and $x$ is a singleton, then $x$
cannot mix with any $t\notin\{x,y\}$. Indeed every nonunit of $P_t$
hits this $t$-free degree, hence contains $x$ or $y$. A jointly
$t,x$-active column would have a $y$-free nonunit hit in $P_t$ by
avoidance, but any hit there is $x$-free by defect-zero vanishing.

If two singletons $a,b$ mix, all three pair supports on the other three
primes occur in both $P_a$ and $P_b$: on the mixed column avoid each
of the other primes in turn. Every singleton among the other three
is therefore isolated by the preceding observation. The same isolation
holds if $n=4620$ and $2$ mixes with an odd singleton $a$. Here is the
extra argument needed at the quartic prime. For a joint $2,a$-active
column $J$ and another odd singleton $w$, choose a $w$-active,
two-inactive column $K$. Such a column occurs in $W_{2'}e_0$:
this vector is two-inactive and multiplication by $W_2$ preserves the
$w$-active subspace and gives the global vacuum column up to phase.
Orthogonality of $J,K$ on $P_2$ forces an odd-degree nonunit common hit.
Indeed $v_2(d_J)\le1$ by activity. For every even-degree $X\in P_2$,
$\chi_J(X)=(d_X/d_J)\chi_X(J)$ vanishes modulo a prime over two:
if $d_J$ is even use $\chi_X(J)\equiv d_J$. The common hit omits
$a,w$, so its degree is the product of the other two odd primes.
Varying $w$ gives all three pair supports and isolates their primes.

An isolated singleton is at least five by Lemma~\ref{lem:isolated-dyadic-square}.
For square-free $n$ the only possible mixed pair is consequently
$2,3$, and it must occur. For $n=4620$ the possible pairs are $2,3$
or $3,a$ with $a=5,7,11$. Indeed $3$ must mix; a pair containing two
isolates the other three primes, and in the absence of such a pair
two is isolated. A mixed odd pair isolates the two remaining odd
primes, so they cannot be partners elsewhere.

We use two numerical consequences of Proposition~\ref{second:two-columns}.
For a set $T$ of isolated singletons, each sign $\epsilon=\pm1$
determines $j_t$ by $1\le j_t<t$ and $(n/t)j_t\equiv-\epsilon\pmod t$.
Necessarily $V_\epsilon=\sum_{t\in T}(n/t)j_t\le n+1$.
In the table, a blank sign entry means that both sums exceed the bound.
\[
\begin{array}{r|c|rr|c@{\qquad}r|c|rr|c}
 n&T&V_+&V_-&\epsilon &n&T&V_+&V_-&\epsilon\\\hline
2310&5,7,11&5004&1926&- &4830&5,7,23&10464&4026&-\\
2730&5,7,13&5004&3186& &4620&5,7,11&5004&8856&\\
3570&5,7,17&4164&6546& &4620&7,11&4080&5160&+\\
3990&5,7,19&7314&4656& &4620&5,11&3024&6216&+\\
4290&5,11,13&5004&7866& &4620&5,7&2904&6336&+
\end{array}
\]

For each surviving row choose $u<v$ in $T$ as below. The only possible
nonunit degree prime to $uv$ is $D=n/(uv)$. To see this, $uv$ occurs on
both principal blocks of the mixed pair, so a degree coprime to $uv$
must contain both mixed primes. If $n$ is square-free it must also
contain the remaining singleton $w$: otherwise it and $uv$ would be
coprime labels in $P_w$. At $n=4620$ the mixed pair is odd and $2$
is isolated. A degree deficient at $2$ lies in $P_2$
(Proposition~\ref{deep:cluster}) and cannot be coprime to $uv$.
This proves the assertion in both cases.

The mixed column identity now requires the products $A_u A_v$ on the
labels of degree $D$ to sum to $-j_uj_v$. At each endpoint $t$
the coordinate is congruent to the inverse of the opposite prime
modulo $t$, and $A_t^2\le j_t(t+1-j_t)$ since $D<n/t$. These two
tests give the following complete comparison.
\[
\begin{array}{r|c|c|c|c}
 n&(u,v)&(j_u,j_v)&(A_u,A_v)& -j_uj_v\\\hline
2310&(5,7)&(3,1)& A_v\equiv3\pmod7,\ A_v^2\le7&-3\\
4830&(5,7)&(1,2)&(-2,3)&-2\\
4620&(7,11)&(3,5)&(2,-3)&-15\\
4620&(5,11)&(1,5)&(1,-2)&-5\\
4620&(5,7)&(1,3)&(-2,3)&-3
\end{array}
\]
The first line admits no $A_v$. In the others the product does not
divide the required sum. This excludes all ordinary square cases.

\paragraph{The reduced square cases.}
If $n$ is square-free, $v_2(n^2/2)=1$ contradicts square-fullness.
It remains that $M=4620^2$. Write $m=1155=3\cdot5\cdot7\cdot11$ and
$\Delta=8m^2$. The positive principal two-block $P$ has odd degrees
and $8\mid |P|$ by Lemma~\ref{lem:fermion-quotients}. A nonunit
of $P$ has neither prime-power degree nor degree the product of two
odd primes by Proposition~\ref{lem:principal-fermion-two}. Thus its
degree is $m$ or $m/a$ for $a\mid m$ prime.

A deep fermion pair exists by the parity assertion of Lemma~\ref{fam:Q:5.9}.
No deep pair is $a$-active for $a\mid m$: its row on $P$ would be
orthogonal to the vacuum and every nonunit hit would omit $a$, hence
would have degree divisible by each other prime of $m$. Reducing
at such a prime leaves only the vacuum.

Lemma~\ref{short:deep-pair} now places every deep pair $Z$ in $P$ with
pure dyadic codegree, hence $f_Z=2$ or $8$. Codegree two would allow at
most one nonunit hit besides the vacuum by the mean-square bound;
dimension orthogonality would then make the inverse of a nonunit degree
integral. Thus $d_Z=m$ and $f_Z=8$. The row hits every member of $P$ by
principality since $m$ is odd. Each nonzero character entry costs at least one
in Galois-averaged squared norm, so $|P|=8$, the row vanishes elsewhere,
and every entry on $P$ is a root of unity by equality and Kronecker's theorem.

Dimension orthogonality modulo $a\mid m$ requires a label of degree
$m/a$ in $P$. If it occurs just once in $P$, Galois invariance of $P$
makes that label rational, so its value on the deep row is $\pm1$.
Thus $m/a\equiv\pm1\pmod a$. But $m/7\equiv4\pmod7$ and
$m/11\equiv6\pmod{11}$. These two degrees occur at least twice.
With the unit and all four required degrees, this leaves exactly
one place in $P$ for a deep label. Thus the deep pair is unique.

The rank-one assertion of Lemma~\ref{short:deep-pair} now gives
$v_2(f_Z)\le2$, contrary to $f_Z=8$. This completes the five-prime case.

\section*{Acknowledgments}
Sebastian Burciu was supported by a grant of the Ministry of Research,
Innovation and Digitization, CNCS--UEFISCDI, project number
PN-IV-P1-PCE-2023-0264, within PNCDI IV.
S\'ebastien Palcoux is supported by the National Natural Science
Foundation of China (NSFC), Grant no.~12471031.

Generative AI assisted with the development and drafting of the arguments and with checking the exposition and arithmetic. All finite checks used in the proof are given in the text.
\begin{thebibliography}{99}\small

\bibitem{BH}
M.~Beck and S.~Ho\c{s}ten, \href{https://doi.org/10.4310/mrl.2006.v13.n4.a10}{\emph{Cyclotomic polytopes and growth series of cyclotomic lattices}}, Math. Res. Lett. \textbf{13} (2006), 607--622.

\bibitem{BGNPWSuper}
P.~Bruillard, C.~Galindo, S.-H.~Ng, J.~Y.~Plavnik, E.~C.~Rowell and Z.~Wang,
\href{https://doi.org/10.1007/s10468-019-09873-9}{\emph{Classification of super-modular categories by rank}},
Algebr. Represent. Theory \textbf{23} (2020), 795--809.

\bibitem{BNRWFinite}
P.~Bruillard, S.-H.~Ng, E.~C.~Rowell, and Z.~Wang,
\href{https://doi.org/10.1090/jams/842}{\emph{Rank-finiteness for modular categories}},
J. Amer. Math. Soc. \textbf{29} (2016), no.~3, 857--881.

\bibitem{BP}
S.~Burciu and S.~Palcoux,
\href{https://doi.org/10.1142/S0219199726500070}{\emph{Burnside-type results for fusion categories}},
Commun. Contemp. Math. (2026), article 2650007.

\bibitem{DLNCongruence}
C.~Dong, X.~Lin, and S.-H.~Ng,
\href{https://doi.org/10.2140/ant.2015.9.2121}{\emph{Congruence property in conformal field theory}},
Algebra \& Number Theory \textbf{9} (2015), no.~9, 2121--2166.

\bibitem{DongWeakIntegral}
J.~Dong,
\href{https://doi.org/10.1007/s11401-026-0019-6}{\emph{Classification of certain weakly integral fusion categories}},
Chinese Ann. Math. Ser. B \textbf{47} (2026), no.~3,
555--572.

\bibitem{DGNO}
V.~Drinfeld, S.~Gelaki, D.~Nikshych, and V.~Ostrik,
\href{https://doi.org/10.1007/s00029-010-0017-z}{\emph{On braided fusion categories I}},
Selecta Math. (N.S.) \textbf{16} (2010), 1--119.

\bibitem{EGNO}
P.~Etingof, S.~Gelaki, D.~Nikshych, and V.~Ostrik,
\href{https://math.mit.edu/~etingof/egnobookfinal.pdf}{\emph{Tensor Categories}},
Mathematical Surveys and Monographs, vol.~205,
American Mathematical Society, 2015.

\bibitem{ENOannals}
P.~Etingof, D.~Nikshych, and V.~Ostrik,
\href{https://doi.org/10.4007/annals.2005.162.581}{\emph{On fusion categories}},
Ann. of Math. \textbf{162} (2005), no.~2, 581--642.

\bibitem{ENOsolvable}
P.~Etingof, D.~Nikshych, and V.~Ostrik,
\href{https://doi.org/10.1016/j.aim.2010.06.009}{\emph{Weakly group-theoretical and solvable fusion categories}},
Adv. Math. \textbf{226} (2011), no.~1, 176--205.

\bibitem{GNNcenters}
S. Gelaki, D. Naidu and D. Nikshych,
\href{https://msp.org/ant/2009/3-8/ant-v3-n8-p05-s.pdf}{\emph{Centers of graded fusion categories}},
Algebra Number Theory \textbf{3} (2009), 959--990.

\bibitem{GNnil}
S. Gelaki and D. Nikshych, \href{https://doi.org/10.1016/j.aim.2007.08.001}{\emph{Nilpotent fusion categories}}, Adv. Math. 217 (2008), 1053--1071.

\bibitem{JFR}
T. Johnson-Freyd and D. Reutter, \href{https://doi.org/10.1090/jams/1023}{\emph{Minimal nondegenerate extensions}}, J. Amer. Math. Soc. \textbf{37} (2024), 81--150.

\bibitem{Malikiosis}
R.~D. Malikiosis, \href{https://doi.org/10.1017/fms.2022.14}{\emph{On the structure of spectral and tiling subsets of cyclic groups}}, Forum Math. Sigma \textbf{10} (2022), e23.

\bibitem{NSTV}
S.~Nasseh, S.~Sather-Wagstaff, R.~Takahashi and K.~VandeBogert,
\href{https://doi.org/10.1016/j.jpaa.2018.06.006}{\emph{Applications and homological properties of local rings with
decomposable maximal ideals}}, J. Pure Appl. Algebra \textbf{223}
(2019), 1272--1287.

\bibitem{NataleWGT}
S.~Natale,
\href{https://doi.org/10.4171/jncg/177}{\emph{On weakly group-theoretical non-degenerate braided fusion categories}},
J. Noncommut. Geom. \textbf{8} (2014), 1043--1060.

\bibitem{NataleCore}
S.~Natale,
\href{https://doi.org/10.1142/S0129167X1850012X}{\emph{The core of a weakly group-theoretical braided fusion category}},
Internat. J. Math. \textbf{29} (2018), no.~2, 1850012.

\bibitem{NSfs}
S.-H. Ng and P. Schauenburg, \href{https://doi.org/10.1016/j.aim.2006.07.017}{\emph{Frobenius--Schur indicators and exponents of spherical categories}}, Adv. Math. \textbf{211} (2007), 34--71.

\bibitem{NSind}
S.-H. Ng and P. Schauenburg, \href{https://doi.org/10.1007/s00220-010-1096-6}{\emph{Congruence subgroups and generalized Frobenius--Schur indicators}}, Comm. Math. Phys. 300 (2010), 1--46.

\bibitem{YuSlightlyDegenerate}
Z.~Yu,
\href{https://doi.org/10.1016/j.jalgebra.2020.04.022}{\emph{On slightly degenerate fusion categories}},
J. Algebra \textbf{559} (2020), 408--431.

\end{thebibliography}

\end{document}
